% =====================================================================
%  TESI DI LAUREA TRIENNALE IN INFORMATICA
%  Universita degli Studi di Camerino
%
%  BOZZA DI LAVORO — generata il 21/08/2026
%  Compilare con pdfLaTeX + Biber (su Overleaf: Menu > Compiler > pdfLaTeX)
% =====================================================================
\documentclass[12pt,a4paper,twoside,openright]{report}

% ---------- Lingua e codifica ----------
\usepackage[T1]{fontenc}
\usepackage[utf8]{inputenc}
\usepackage[italian]{babel}
\usepackage{csquotes}

% ---------- Layout ----------
\usepackage[a4paper,inner=3.5cm,outer=2.5cm,top=3cm,bottom=3cm]{geometry}
\usepackage{setspace}
\onehalfspacing
\usepackage{emptypage}
\usepackage{fancyhdr}
% fancyhdr chiede almeno 14.5pt di testatina: senza questa riga avvisa su ogni pagina
\setlength{\headheight}{14.5pt}
\addtolength{\topmargin}{-2.5pt}
\pagestyle{fancy}
\fancyhf{}
\fancyhead[LE,RO]{\thepage}
\fancyhead[RE]{\nouppercase{\leftmark}}
\fancyhead[LO]{\nouppercase{\rightmark}}
\renewcommand{\headrulewidth}{0.4pt}

% ---------- Matematica, tabelle, figure ----------
\usepackage{amsmath,amssymb}
\usepackage{booktabs}
\usepackage{longtable}
\usepackage{multirow}
\usepackage{tabularx}
\usepackage{array}
\usepackage{graphicx}
\usepackage{float}
\usepackage[font=small,labelfont=bf]{caption}
\usepackage{subcaption}
\usepackage{xcolor}
\usepackage{enumitem}
\usepackage{textcomp}
\usepackage{gensymb}   % \degree, \celsius
\usepackage{eurosym}   % \euro

% ---------- Codice ----------
\usepackage{listings}
\definecolor{codegray}{gray}{0.95}
\definecolor{codegreen}{rgb}{0,0.5,0}
\definecolor{codeblue}{rgb}{0,0,0.6}
\lstset{
  backgroundcolor=\color{codegray},
  basicstyle=\ttfamily\footnotesize,
  keywordstyle=\color{codeblue}\bfseries,
  commentstyle=\color{codegreen}\itshape,
  stringstyle=\color{red!60!black},
  breaklines=true,
  showstringspaces=false,
  frame=single,
  framesep=4pt,
  numbers=left,
  numberstyle=\tiny\color{gray},
  captionpos=b
}
% Nota: i listati sono volutamente in ASCII puro. Se serve inserire lettere
% accentate dentro un lstlisting, aggiungere qui la riga:
%   literate={a'}{{\`a}}1 {e'}{{\`e}}1 ...
% oppure usare extendedchars=true con l'opportuna codifica.

% ---------- Bibliografia ----------
\usepackage[backend=biber,style=numeric,sorting=none,maxbibnames=6]{biblatex}
\addbibresource{bib/tesi.bib}

% ---------- Hyperref (per ultimo) ----------
\usepackage[hidelinks,breaklinks]{hyperref}
\usepackage{cleveref}

% ---------- Comandi di comodo ----------
\newcommand{\ld}{LD2410B}
\newcommand{\ldb}{HLK-LD2410B}
\newcommand{\ldc}{HLK-LD2420}
\newcommand{\pir}{HC-SR501}
\newcommand{\mmw}{mmWave}
% Nota redazionale visibile solo in bozza: commentare la riga sotto per la stampa finale
\newcommand{\todo}[1]{\textcolor{red}{\textbf{[TODO: #1]}}}
% \renewcommand{\todo}[1]{}

\begin{document}

% La classe "report" non definisce \frontmatter/\mainmatter:
% la numerazione delle pagine e' gestita a mano qui sotto.
\pagenumbering{roman}

% La pagina delle proposte di titolo va rimossa nella versione finale
% =====================================================================
%  PAGINA DI LAVORO — proposte di titolo
%  Da eliminare (o commentare la riga in main.tex) nella versione consegnata.
% =====================================================================
\thispagestyle{empty}

\begin{center}
{\Large\bfseries Proposte di titolo}\\[2pt]
{\small documento di lavoro --- 21 agosto 2026}
\end{center}

\vspace{0.5em}
\noindent
Elenco di titoli alternativi, raggruppati per taglio. La scelta va concordata con
il relatore: il titolo definitivo va poi riportato nel file del frontespizio
(\texttt{capitoli/}\allowbreak\texttt{00-frontespizio.tex}) e questa pagina va
eliminata.

\vspace{1em}
\noindent\textbf{A. Taglio ``confronto sperimentale'' (il più aderente al contenuto)}
\begin{enumerate}[label=A\arabic*.,leftmargin=2.2em,itemsep=3pt]
  \item Sensori di presenza per la sicurezza sismica: confronto sperimentale tra
        radar mmWave a 24~GHz e sensori PIR su piattaforma ESP32
  \item Rilevare chi non si muove: valutazione sperimentale di radar mmWave e PIR
        per il rilevamento di presenza in scenari di emergenza
  \item Caratterizzazione e confronto quantitativo di sensori PIR e radar mmWave a
        24~GHz per il rilevamento di persone in ambienti indoor
  \item Oltre il movimento: analisi comparata PIR / mmWave FMCW per il rilevamento
        di presenza statica
\end{enumerate}

\vspace{0.6em}
\noindent\textbf{B. Taglio ``progetto UPRISE/SAFE'' (contesto applicativo in evidenza)}
\begin{enumerate}[label=B\arabic*.,leftmargin=2.2em,itemsep=3pt]
  \item Sensing di presenza per arredi salva-vita: il radar mmWave a 24~GHz come
        evoluzione del nodo PIR nel progetto SAFE
  \item Un sensore per chi è rimasto sotto: rilevamento di presenza e di vitalità
        con radar mmWave per il soccorso post-sismico
  \item Dal ``c'è qualcuno'' al ``in che condizioni è'': sensing mmWave a supporto
        del triage in emergenza sismica
  \item Nodi IoT per arredi antisismici: valutazione di sensori mmWave, PIR e UWB
        per il rilevamento di persone intrappolate
\end{enumerate}

\vspace{0.6em}
\noindent\textbf{C. Taglio ``sistema realizzato'' (firmware, web UI, indice di vitalità)}
\begin{enumerate}[label=C\arabic*.,leftmargin=2.2em,itemsep=3pt]
  \item Progetto e realizzazione di un nodo di sensing mmWave su ESP32 con
        dashboard web e indice di vitalità
  \item Acquisizione, visualizzazione e analisi di dati radar mmWave a 24~GHz: un
        sistema completo su ESP32 per il rilevamento di presenza
  \item Da micro-movimenti a indice di vitalità: elaborazione di dati radar mmWave
        su microcontrollore
\end{enumerate}

\vspace{0.6em}
\noindent\textbf{D. Taglio tecnico e breve}
\begin{enumerate}[label=D\arabic*.,leftmargin=2.2em,itemsep=3pt]
  \item Radar mmWave a 24~GHz per il rilevamento di presenza statica: misure e
        confronto con sensori PIR
  \item Presence sensing con mmWave ed ESP32: caratterizzazione sperimentale
  \item Micro-movimenti e respiro: analisi dei dati per-gate di un radar FMCW a
        24~GHz
\end{enumerate}

\vspace{1.2em}
\noindent\textbf{Nota per la scelta.} Il risultato sperimentale centrale già
acquisito è la \emph{cecità del PIR sulla persona ferma}: falsi negativi al
$99{,}78 \pm 0{,}49\,\%$ contro lo $0{,}00\,\%$ del radar
(\cref{sec:persona-ferma}). I titoli A2, B2 e D1 sono quelli che lo mettono in
prima linea. Se il relatore preferisce che il titolo espliciti il progetto europeo,
la coppia consigliata è B1 (sobria) oppure B2 (più comunicativa). Il taglio C è
appropriato solo se l'implementazione della web UI e dell'indice di vitalità
saranno completate: al momento sono progettate ma non realizzate.

\cleardoublepage

% =====================================================================
%  FRONTESPIZIO
%  Sostituire il titolo con quello scelto fra le proposte della pagina
%  precedente e verificare la dicitura ufficiale del corso di laurea.
% =====================================================================
\begin{titlepage}
\centering
\thispagestyle{empty}

{\large UNIVERSIT\`A DEGLI STUDI DI CAMERINO\par}
\vspace{4pt}
{\large Scuola di Scienze e Tecnologie\par}
\vspace{2pt}
{\normalsize Corso di Laurea in Informatica (L-31)\par}

\vspace{1.2cm}
% \includegraphics[width=0.28\textwidth]{img/logo_unicam.pdf}   % <-- inserire il logo
\vspace{1.2cm}

{\Large\bfseries
Sensori di presenza per la sicurezza sismica:\par
confronto sperimentale tra radar mmWave a 24~GHz\par
e sensori PIR su piattaforma ESP32\par}

\vspace{0.4cm}
{\small\itshape (titolo provvisorio --- vedi pagina delle proposte)\par}

\vspace{2.2cm}

\begin{minipage}[t]{0.46\textwidth}
\raggedright
\textbf{Relatore}

\vspace{2pt}
Dott. Massimo Callisto De Donato
\end{minipage}%
\hfill
\begin{minipage}[t]{0.46\textwidth}
\raggedleft
\textbf{Laureando}

\vspace{2pt}
Stefano De Bernardinis
\end{minipage}

\vfill
{\normalsize Anno Accademico 2025/2026\par}
\end{titlepage}
\cleardoublepage


\cleardoublepage
\chapter*{Sommario}
\addcontentsline{toc}{chapter}{Sommario}
\markboth{Sommario}{Sommario}

Il progetto europeo UPRISE, nato dopo il sisma del 2016 nell'area di Camerino,
prevede l'integrazione di sensori IoT negli arredi scolastici: in caso di terremoto
il nodo installato nell'arredo deve rilevare la presenza di una persona rifugiata
sotto di esso e comunicarlo all'esterno tramite rete LoRa. Il nodo attualmente in
uso (DIPME-DEVICE) affida il rilevamento di presenza a un sensore \emph{PIR}, che
per principio fisico è cieco verso una persona immobile: proprio la condizione
tipica di chi si trova intrappolato.

Questa tesi affronta la parte di \emph{sensing} del problema. Sono stati studiati,
strumentati e confrontati sperimentalmente due radar \emph{mmWave} FMCW a 24~GHz
(HLK-LD2410B e HLK-LD2420) e un sensore PIR HC-SR501, acquisendo i dati con una
piattaforma di misura realizzata su ESP32 che campiona a 5~Hz in
\emph{engineering mode}, cioè con accesso ai valori di energia riflessa per ciascuno
dei nove \emph{gate} di distanza del radar.

Il risultato centrale è quantitativo: su cinque prove indipendenti da cinque minuti
ciascuna, con un soggetto immobile seduto a 2,30~m, il PIR ha mancato la presenza
nel $99{,}78 \pm 0{,}49\,\%$ dei campioni, mentre il radar l'ha rilevata nel
$100\,\%$ dei 7554 campioni acquisiti. Il PIR risulta inaffidabile anche con il
soggetto \emph{in movimento sul posto} ($80{,}5\,\%$ di falsi negativi a 1~m,
$100\,\%$ oltre i 2~m), coerentemente con il fatto che l'elemento piroelettrico
risponde al transito del flusso infrarosso fra le zone della lente di Fresnel e non
al movimento in sé. La misura di distanza del radar si è rivelata estremamente
lineare ($R^2 = 0{,}99965$ su 25 prove fra 1 e 5~m), con un errore di scala del
$+3{,}81\,\%$ correggibile mediante un singolo coefficiente; su 7~ore di stanza
vuota nessuno dei due sensori ha prodotto falsi positivi, con limite superiore al
95\,\% di confidenza pari a $0{,}43$ eventi/h. Dall'analisi in frequenza
dell'energia per-gate è stato inoltre possibile estrarre il ritmo respiratorio di un
soggetto immobile, con concordanza fra canali indipendenti.

Su questa base la tesi propone un \emph{indice di vitalità}: una grandezza compresa
fra 0 e 100 che, oltre al ``c'è qualcuno'', stimi quanto la persona si muove, a
supporto del triage dei soccorritori. Il lavoro comprende infine l'analisi dei
consumi da datasheet e il progetto di un'interfaccia web servita direttamente
dall'ESP32 per la visualizzazione in tempo reale e l'esportazione dei dati.

La raccomandazione conclusiva non è la sostituzione del PIR, ma
un'\emph{architettura ibrida}: il PIR come guardiano a microampere che decide quando
accendere il radar, e il radar mmWave come strumento di misura durante la finestra
di emergenza.

\vspace{1em}
\begin{sloppypar}
\noindent\textbf{Parole chiave:} radar mmWave, FMCW 24~GHz, PIR, rilevamento di
presenza, ESP32, HLK-LD2410B, monitoraggio del respiro, sicurezza sismica,
UPRISE/SAFE.
\end{sloppypar}


\cleardoublepage
\tableofcontents
\cleardoublepage
\listoffigures
\cleardoublepage
\listoftables

\cleardoublepage
\pagenumbering{arabic}

\chapter{Introduzione}
\label{cap:introduzione}

\section{Contesto: il progetto UPRISE/SAFE}
\label{sec:contesto}

Il terremoto del 2016 che ha colpito l'Italia centrale, e in particolare l'area di
Camerino, ha reso evidente un problema tanto concreto quanto poco studiato: nelle
prime ore successive a un evento sismico l'informazione più preziosa per i
soccorritori non è lo stato dell'edificio, ma \emph{dove si trovano le persone}.
Da questa esigenza nasce il progetto europeo UPRISE\footnote{Nella documentazione
scientifica del gruppo di ricerca il progetto è indicato anche con il nome
\textbf{SAFE}; nel seguito si userà UPRISE/SAFE per riferirsi allo stesso
ecosistema.}, che ha come obiettivo la mitigazione del rischio sismico attraverso
l'integrazione di sensoristica IoT negli \emph{arredi} degli ambienti
scolastici~\cite{slideUprise}.

L'idea del progetto si articola su tre livelli.

\begin{description}[leftmargin=1.4em,style=nextline]
  \item[Arredi salva-vita] Banchi e pareti progettate dalla Scuola di Architettura
    e Design di UNICAM per resistere al crollo e offrire un volume di rifugio: il
    banco ha telaio doppio spaziale, piano rinforzato con lamiera forata
    antisfondamento, elemento di base dissipativo e connessione strutturale fra
    banchi adiacenti.
  \item[Nodi di sensing] Nell'arredo è integrata l'elettronica sviluppata dal
    partner industriale AM Microsystems: il \textbf{DIPME-DEVICE} monta radio LoRa,
    sensore \emph{PIR}, sensore \emph{UWB}, oltre a CO\textsubscript{2},
    temperatura, umidità, pressione, accelerometro e Bluetooth. I nodi
    comunicano con un \textbf{DIPME-COORDINATOR} e da questo verso un gateway IoT
    su Linux/Raspberry.
  \item[Raccolta e visualizzazione] In emergenza la rete di terra non è
    disponibile: i messaggi LoRa (bande sub-GHz libere, 433/868~MHz in Europa,
    con portate dell'ordine di 3--5~km in ambiente urbano) vengono raccolti da
    \emph{droni} che ricostruiscono una mappa di calore geolocalizzata, consultabile
    offline dal tablet dei soccorritori. L'intero sistema è modellato e validato
    a priori sulla piattaforma di Digital Twin del gruppo di
    ricerca~\cite{callisto2023dt,unicamDT}.
\end{description}

Il sistema opera in due modalità: un ``tempo di pace'' di monitoraggio ordinario
e un ``tempo di guerra'' di emergenza sismica, attivato da un trigger (accelerometro
sul gateway, blackout della rete elettrica).

\todo{inserire una figura di sistema: arredo $\to$ DIPME $\to$ coordinator $\to$
gateway $\to$ drone $\to$ tablet}

\section{Il problema specifico affrontato dalla tesi}
\label{sec:problema}

Questa tesi si occupa esclusivamente del livello più basso della catena: il
\textbf{sensing di presenza}. La domanda di partenza è molto circoscritta e ha
un'origine precisa nell'architettura esistente.

Il DIPME-DEVICE rileva la presenza con un sensore PIR (\emph{Passive InfraRed}).
Il PIR è il sensore di presenza più diffuso al mondo per tre ragioni difficili
da battere: consuma qualche decina di microampere, costa pochi euro e non richiede
alcuna elaborazione. Ha tuttavia un limite che non dipende dalla qualità del
componente ma dal suo principio di funzionamento: l'elemento piroelettrico risponde
alla \emph{variazione} del flusso infrarosso incidente, non al suo valore. Una
persona ferma produce un flusso costante, quindi un segnale differenziale nullo:
\emph{nessun PIR, a nessun prezzo, può rilevare una persona immobile}
(\cref{sec:pir-principio}).

Nello scenario UPRISE questo limite cade esattamente nel punto peggiore. Una
persona rifugiata sotto un banco dopo un crollo è, tipicamente, \textbf{ferma}:
per paura, per scelta, o perché non può muoversi. È il caso cieco del sensore
che il nodo monta oggi.

I radar \emph{mmWave} a onda continua modulata in frequenza (FMCW) a 24~GHz
promettono di coprire proprio questo buco: rilevano i micro-movimenti
(fino, in linea di principio, all'oscillazione toracica del respiro), misurano la
distanza del bersaglio e non producono immagini, il che li rende compatibili con i
vincoli di privacy di un ambiente scolastico. Da qui le domande di ricerca della
tesi:

\begin{enumerate}[leftmargin=2em,itemsep=2pt]
  \item \emph{quanto} vale, in numeri, il vantaggio del mmWave sul PIR nel caso
        della persona ferma?
  \item quali \emph{dati} produce realmente un modulo mmWave di fascia consumer, e
        con quale accuratezza?
  \item a che \emph{costo} energetico, e con quale architettura di nodo?
  \item è possibile andare oltre la presenza binaria, verso una misura di
        \emph{vitalità} utile al triage?
\end{enumerate}

\section{Obiettivi}
\label{sec:obiettivi}

Gli obiettivi sono stati concordati con il relatore e sono qui riportati nella
formulazione originale, perché costituiscono la traccia di tutto il documento.

\begin{enumerate}[label=\textbf{O\arabic*.},leftmargin=2.6em,itemsep=4pt]
  \item \textbf{Studio dei sensori} --- chi li usa, come vengono usati, quali dati
        producono (PIR e mmWave). \emph{Capitolo \ref{cap:sensori}.}
  \item \textbf{Analisi dei dati forniti} --- capire nel dettaglio quali grandezze
        i sensori mettono a disposizione: il PIR solo movimento sì/no, il mmWave
        micro-movimenti, distanza e potenzialmente respiro.
        \emph{Capitolo \ref{cap:moduli}.}
  \item \textbf{Comparazione con testing numerico} --- confronto sperimentale
        PIR vs mmWave con metriche quantitative. \emph{Capitolo \ref{cap:confronto}.}
  \item \textbf{Consumo energetico} --- a titolo informativo, sulla base dei
        datasheet, non essendo disponibile strumentazione di misura.
        \emph{Capitolo \ref{cap:consumi}.}
  \item \textbf{User interface web} --- l'ESP32 pubblica i dati, un sito web li
        visualizza in tempo reale, la stessa applicazione li salva con statistiche
        ed esporta in CSV per l'analisi numerica. \emph{Capitolo \ref{cap:webui}.}
  \item \textbf{Indice di vitalità} --- obiettivo avanzato: oltre alla presenza
        sì/no, un algoritmo che quantifichi quanto una persona si muove
        (per esempio 50 su 100) e la classifichi di conseguenza.
        \emph{Capitolo \ref{cap:vitalita}.}
\end{enumerate}

\paragraph{Fuori scope.} Non fanno parte della tesi: la radio LoRa e il suo
protocollo, la raccolta con droni, il Digital Twin, la progettazione strutturale
degli arredi. Questi elementi compaiono solo in \cref{sec:contesto} come
inquadramento e, dove pertinente, come vincolo di progetto (per esempio l'assenza
di connettività in emergenza, che orienta le scelte del \cref{cap:webui}).

\section{Contributi}
\label{sec:contributi}

I contributi originali di questo lavoro sono:

\begin{itemize}[leftmargin=1.4em,itemsep=3pt]
  \item una \textbf{piattaforma di misura riproducibile} su ESP32 che acquisisce
        radar e PIR in parallelo a 5~Hz esatti con accesso ai dati per-gate del
        radar, con formato CSV unico condiviso da tutta la catena di analisi
        (\cref{sec:piattaforma});
  \item la \textbf{caratterizzazione quantitativa} del divario PIR/mmWave sulla
        persona ferma e sulla persona in movimento sul posto, con deviazioni
        standard su cinque ripetizioni per scenario (\cref{cap:confronto});
  \item la \textbf{calibrazione della misura di distanza} del LD2410B su 25 prove
        fra 1 e 5~m, con individuazione di un errore di scala sistematico e del
        coefficiente che lo corregge (\cref{sec:distanza});
  \item la \textbf{dimostrazione di misurabilità del respiro} con hardware
        consumer, e l'individuazione documentata di quale dei venti canali di
        energia disponibili sia effettivamente utilizzabile a ciascuna distanza
        (\cref{sec:respiro});
  \item la \textbf{specifica e la taratura di un indice di vitalità}
        implementabile su microcontrollore (\cref{cap:vitalita}).
\end{itemize}

\section{Struttura del documento}

Il \cref{cap:sensori} tratta i principi di funzionamento e lo stato dell'arte dei
sensori di presenza (O1). Il \cref{cap:moduli} descrive i moduli in esame, i dati
che producono e la piattaforma di acquisizione realizzata (O2). Il
\cref{cap:confronto} è il capitolo centrale e riporta il confronto sperimentale
(O3). Il \cref{cap:consumi} analizza i consumi (O4), il \cref{cap:webui} la web UI
(O5), il \cref{cap:vitalita} l'indice di vitalità (O6). Il
\cref{cap:conclusioni} risponde in sintesi alla domanda di partenza e indica gli
sviluppi futuri. Le appendici raccolgono i dettagli del protocollo seriale e i
listati principali.

\chapter{I sensori di presenza: principi e stato dell'arte}
\label{cap:sensori}

\begin{quote}\small
Questo capitolo risponde all'obiettivo O1: che cosa sono i sensori di presenza,
come funzionano, chi li usa e quali dati producono. La trattazione è
volutamente asimmetrica --- più spazio a PIR e mmWave, che sono i sensori
effettivamente misurati nel \cref{cap:confronto} --- e si chiude con il
posizionamento comparativo che giustifica le scelte del resto della tesi.
\end{quote}

\section{Il sensore PIR}
\label{sec:pir}

\subsection{Principio fisico: un rilevatore di variazione, non di presenza}
\label{sec:pir-principio}

PIR è l'acronimo di \emph{Passive InfraRed}: il sensore non emette nulla, si limita
a osservare la radiazione infrarossa che lo raggiunge. Il cuore del componente è un
cristallo \textbf{piroelettrico}, che genera una piccola carica elettrica quando la
radiazione incidente \emph{cambia}. Il corpo umano, con una temperatura superficiale
di circa 34\,\celsius, emette con picco intorno a 9,4~\textmu m; il sensore è
filtrato con una finestra ottica di 5--14~\textmu m proprio per privilegiare quella
banda~\cite{adafruitPIR}.

Il punto essenziale per questa tesi è che il cristallo risponde alla
\textbf{derivata} del flusso infrarosso, non al suo valore assoluto.
Costruttivamente il sensore contiene \textbf{due elementi piroelettrici in
configurazione differenziale}, collegati in serie con polarità opposte, e da questo
seguono tre conseguenze:

\begin{itemize}[leftmargin=1.4em,itemsep=2pt]
  \item una persona che si muove investe i due elementi in \emph{tempi diversi}: il
        segnale differenziale è non nullo e si ha rilevamento;
  \item una persona ferma illumina entrambi gli elementi con lo stesso flusso
        costante: la differenza è nulla e il sensore è \textbf{fisicamente cieco};
  \item una variazione globale della scena (il sole che scalda la stanza) colpisce
        entrambi gli elementi allo stesso modo e si cancella: da qui la buona
        immunità del PIR ai falsi positivi ambientali lenti.
\end{itemize}

È importante inquadrare correttamente questo comportamento. Il risultato
sperimentale del \cref{sec:persona-ferma} --- falsi negativi prossimi al
100\,\% su persona ferma --- \textbf{non è il fallimento di un componente
economico}: è il funzionamento nominale del principio. Nessun PIR, a nessun prezzo,
può rilevare un corpo immobile. Questa è la frase su cui poggia l'intero confronto
con il radar.

\subsection{La lente di Fresnel: da due elementi a decine di zone}
\label{sec:fresnel}

Due soli elementi coprirebbero un campo visivo strettissimo. La cupola bianca che
si vede su tutti i moduli PIR è una \textbf{lente di Fresnel multi-segmento}: ogni
segmento focalizza sul cristallo una porzione differente della scena, creando
decine di \emph{zone di rilevamento} a ventaglio, separate da zone cieche.

Questo dettaglio costruttivo ha due conseguenze misurabili, entrambe verificate
sperimentalmente nel \cref{cap:confronto}:

\begin{enumerate}[leftmargin=1.8em,itemsep=2pt]
  \item il movimento \textbf{trasversale}, che attraversa le zone, genera molte
        transizioni di flusso ed è rilevato bene; il movimento \textbf{radiale},
        diretto verso il sensore, resta più a lungo nella stessa zona ed è rilevato
        male. Il radar FMCW ha la sensibilità esattamente opposta, perché misura
        variazioni di distanza;
  \item la portata dichiarata (3--7~m) vale per il movimento trasversale di tutto
        il corpo. Un movimento \emph{sul posto} non attraversa le zone e cade nel
        caso sfavorevole: è il meccanismo che spiega il risultato, a prima vista
        sorprendente, del \cref{sec:pir-movimento}.
\end{enumerate}

\subsection{I dati prodotti: un bit, e come va interpretato}
\label{sec:pir-dati}

L'uscita di un modulo PIR è \textbf{un solo bit digitale}. Tutta l'elaborazione è
nell'elettronica di condizionamento, tipicamente il circuito integrato
BISS0001~\cite{biss0001}, che introduce i meccanismi riassunti nella
\cref{tab:pir-meccanismi}.

\begin{table}[htbp]
\centering
\caption{Meccanismi di condizionamento del segnale PIR e loro effetto sul dato
prodotto.}
\label{tab:pir-meccanismi}
\small
\begin{tabularx}{\textwidth}{@{}lX@{}}
\toprule
\textbf{Meccanismo} & \textbf{Effetto sul dato} \\
\midrule
Soglia sul segnale amplificato & sensibilità/portata (trimmer) \\
\textbf{Tempo di ritenuta} (hold) & l'uscita resta alta $N$ secondi dopo l'ultimo
  trigger (trimmer) \\
Modalità trigger L (single) & l'uscita torna bassa alla fine della ritenuta anche
  se il movimento continua \\
Modalità trigger H (repeat) & ogni nuovo movimento riarma la ritenuta \\
Tempo di stabilizzazione & 30--60~s all'accensione, uscita non attendibile \\
Block time & 2,5~s di cecità strutturale dopo ogni rilascio \\
\bottomrule
\end{tabularx}
\end{table}

Da questa tabella derivano direttamente tre scelte del protocollo sperimentale
(\cref{sec:protocollo}): il jumper in modalità \textbf{H}, il trimmer di ritenuta
al \textbf{minimo} e lo scarto del primo minuto di ogni acquisizione. Se la
ritenuta non fosse al minimo, la ``latenza di rilascio'' misurata sarebbe una
proprietà del potenziometro, non del sensore.

Il confronto informativo con il mmWave si riassume in una riga: il PIR produce
\textbf{1 bit filtrato da isteresi}, il LD2410B produce circa \textbf{venti valori
numerici} (stato, due distanze, due energie di target, diciotto energie per-gate) a
ogni frame. La tesi quantifica quanto questo divario informativo si traduca in
capacità di rilevamento.

\subsection{Sensibilità alla temperatura ambiente}
\label{sec:pir-temperatura}

Il segnale utile del PIR è il \emph{contrasto termico} fra corpo e sfondo. A
20\,\celsius\ di ambiente il contrasto è di circa 14\,\celsius; a
30--32\,\celsius\ --- un'aula assolata d'estate --- scende a 2--4\,\celsius, con
segnale più debole, portata ridotta e più falsi negativi. Sopra i 35\,\celsius\ il
contrasto può addirittura invertirsi.

Per questo motivo ogni sessione sperimentale registra la temperatura della stanza
(\cref{sec:registro}). Le acquisizioni di questa tesi sono state svolte fra 25 e
29\,\celsius: una condizione \emph{sfavorevole} al PIR, che va dichiarata
esplicitamente nella discussione dei risultati (\cref{sec:limiti}).

Altre sorgenti note di falsi positivi PIR, tenute fuori dal setup: correnti d'aria
calda o fredda, animali, sole diretto attraverso i vetri (il vetro blocca l'IR del
corpo, ma non il riscaldamento solare degli oggetti interni).

\subsection{Chi lo usa e perché domina}
\label{sec:pir-usi}

Il PIR è il sensore di presenza più diffuso al mondo: luci automatiche, antifurti,
videocitofoni, termostati. Le ragioni sono tre numeri difficili da battere:
circa \textbf{50~\textmu A} di corrente di riposo, \textbf{1--3~\euro} di costo e
\textbf{zero elaborazione} richiesta al microcontrollore, cui si aggiunge l'assenza
totale di emissioni --- essendo un sensore passivo --- che semplifica
certificazioni e conformità in materia di privacy.

Nel progetto UPRISE il DIPME-DEVICE monta un PIR esattamente per questi motivi: un
nodo deve poter vegliare per anni a batteria. Il limite emerge solo nello scenario
di emergenza, ed è il punto di partenza di questa tesi.

\section{Il radar mmWave a 24~GHz}
\label{sec:mmwave}

\subsection{Principio: FMCW e gate di distanza}

I moduli in esame sono radar \emph{FMCW} (\emph{Frequency Modulated Continuous
Wave}) operanti nella banda ISM a 24~GHz. Il principio è il seguente: il
trasmettitore emette un segnale continuo la cui frequenza cresce linearmente nel
tempo (\emph{chirp}); l'eco riflessa dal bersaglio torna con un ritardo
proporzionale alla distanza, e poiché la frequenza istantanea è cambiata nel
frattempo, il battimento fra segnale trasmesso e ricevuto ha una frequenza
proporzionale alla distanza del bersaglio. Un'analisi in frequenza del segnale di
battimento produce quindi direttamente un profilo distanza--intensità.

I moduli consumer non espongono il segnale grezzo, ma una sua versione elaborata e
discretizzata in \textbf{gate di distanza}: intervalli contigui di profondità, nove
per il LD2410B (0,75~m ciascuno, per un massimo di 6,75~m) e quindici per il
LD2420 (0,70~m ciascuno). Per ogni gate il modulo riporta un valore di
\textbf{energia riflessa} normalizzato nell'intervallo 0--100, separatamente per la
componente ``in movimento'' e per quella ``stazionaria''. Questo insieme di valori
--- accessibile solo attivando l'\emph{engineering mode}, \cref{sec:engineering} ---
è la materia prima di tutta l'analisi quantitativa di questa tesi.

\subsection{Come si rileva un corpo immobile}

Un radar FMCW puro misura distanza, non movimento: un corpo fermo produce comunque
un'eco. In pratica i moduli commerciali separano due contributi:

\begin{itemize}[leftmargin=1.4em,itemsep=2pt]
  \item la componente \textbf{moving}, associata a variazioni rapide della fase
        dell'eco (spostamento del bersaglio);
  \item la componente \textbf{stationary} (o \emph{motionless}), associata a
        un'eco stabile in distanza ma con fluttuazioni residue.
\end{itemize}

Una persona ferma non è mai realmente ferma: il torace si muove di alcuni
millimetri a ogni atto respiratorio, e a 24~GHz (lunghezza d'onda di circa
12,5~mm) uno spostamento millimetrico produce una variazione di fase perfettamente
misurabile. È questa la ragione fisica per cui il radar vede ciò che il PIR non può
vedere, ed è anche la base del rilevamento del respiro
(\cref{sec:respiro-teoria}).

\subsection{Il respiro come segnale}
\label{sec:respiro-teoria}

Il respiro di un adulto a riposo ha frequenza compresa fra circa 0,1 e 0,5~Hz
(6--30 atti al minuto). L'oscillazione toracica modula l'energia riflessa dal gate
in cui si trova la persona: campionando quel valore nel tempo e applicando una
trasformata di Fourier, il picco spettrale in banda 0,1--0,5~Hz corrisponde al
ritmo respiratorio.

Questa strada è ampiamente documentata in letteratura, con hardware sia
professionale sia consumer. Sono stati esaminati in particolare i lavori
sul monitoraggio non a contatto dei parametri vitali con radar FMCW, che
dimostrano il rilevamento di eventi apnoici e confrontano architetture a onda
continua e FMCW~\cite{upm2024vitals}, l'estrazione congiunta di frequenza
respiratoria e cardiaca con reti neurali~\cite{uts2024vitals}, il confronto fra
sensori a 24 e a 60~GHz (dove il 24~GHz privilegia la portata e il 60~GHz la
risoluzione)~\cite{theseus2024}, il tracciamento multi-bersaglio~\cite{uts2023mot}
e le tecniche specifiche di monitoraggio respiratorio con
mmWave~\cite{mdpi2024mmwaverm}.

L'aspettativa realistica con hardware consumer, che i risultati del
\cref{sec:respiro} confermano, è la seguente: rilevare la \emph{presenza} di
respiro --- cioè distinguere una persona ferma e viva da una stanza vuota --- è
affidabile; misurare la frequenza esatta richiede post-elaborazione e condizioni
controllate.

\subsection{Chi li usa}

I radar mmWave di fascia consumer si sono diffusi rapidamente nella domotica
(interruttori di presenza, illuminazione che non si spegne mentre si legge fermi
sul divano), nell'automazione degli edifici per il conteggio di occupazione, nel
monitoraggio del sonno e nei sistemi di rilevamento cadute per l'assistenza agli
anziani. Il supporto nelle piattaforme aperte è ormai maturo: le integrazioni
ESPHome per LD2410 e LD2420 sono documentate e ampiamente
usate~\cite{esphomeLD2410,esphomeLD2420,espboards}, e la letteratura tecnica sulle
applicazioni di conteggio dell'occupazione negli edifici
intelligenti~\cite{mmwaveOccupancy} documenta gli ordini di grandezza attesi. Sono
fonti utili per il
comportamento pratico dei moduli oltre a quanto scritto nei manuali.

\section{Altre tecnologie di rilevamento di presenza}
\label{sec:altre-tecnologie}

\subsection{UWB (Ultra-WideBand)}
\label{sec:uwb}

L'UWB merita più attenzione delle altre alternative per una ragione precisa: il
DIPME-DEVICE ne monta già uno. La domanda ``perché mmWave e non UWB?'' è quindi
attesa in sede di discussione, e la risposta va preparata.

L'UWB trasmette impulsi molto brevi su una banda larghissima (tipicamente
3--10~GHz), ottenendo un'eccellente risoluzione temporale e quindi spaziale, buona
penetrazione dei materiali e consumi molto contenuti (sotto i 10~mW in
trasmissione)~\cite{rfwireless}. È in grado di rilevare corpi fermi e, in
letteratura, anche parametri vitali.

Gli argomenti a favore del mmWave nel contesto specifico di questa tesi sono
quattro:

\begin{enumerate}[leftmargin=1.8em,itemsep=2pt]
  \item \textbf{maturità dei moduli consumer}: LD2410B e LD2420 costano 5--15~\euro,
        hanno firmware che espone direttamente presenza, distanza ed energia
        per-gate e non richiedono elaborazione a valle; i moduli UWB comparabili
        costano 20--50~\euro\ e tipicamente espongono dati grezzi che richiedono un
        processore vero;
  \item \textbf{dati per-gate già elaborati}: i venti canali di energia del
        LD2410B sono ciò che rende possibile l'indice di vitalità del
        \cref{cap:vitalita} su un microcontrollore da pochi euro;
  \item \textbf{destinazione d'uso}: l'UWB sul DIPME è pensato principalmente per
        localizzazione/ranging fra nodi, un problema diverso dal rilevamento di
        presenza in un volume;
  \item \textbf{complementarità, non concorrenza}: nulla vieta che il nodo finale
        usi entrambi. La raccomandazione del \cref{cap:conclusioni} è
        un'architettura a più sensori, non una sostituzione.
\end{enumerate}

\subsection{Ultrasuoni, CO\textsubscript{2}, visione}

Per completezza si citano tre alternative escluse a priori.
Gli \textbf{ultrasuoni} misurano distanza in modo economico ma sono molto
sensibili a tessuti fonoassorbenti e alla geometria della scena, e sotto un arredo
crollato la propagazione acustica è imprevedibile.
La \textbf{CO\textsubscript{2}} è un ottimo indicatore di occupazione di un
ambiente --- ed è già presente sul DIPME --- ma ha costanti di tempo dell'ordine
dei minuti e non localizza: inadatta a rispondere alla domanda ``c'è qualcuno
sotto \emph{questo} banco?''.
La \textbf{visione} (telecamere, termocamere) è la più informativa ma è esclusa da
vincoli di privacy in ambiente scolastico, dal costo e dall'inutilizzabilità nel
buio e nella polvere post-crollo.

\section{Posizionamento comparativo}
\label{sec:posizionamento}

La \cref{tab:comparativa} riassume il confronto fra le tre tecnologie candidate. I
valori sono da datasheet o da letteratura; la colonna mmWave è quella verificata
sperimentalmente nel \cref{cap:confronto}.

\begin{table}[htbp]
\centering
\caption{Confronto fra le tecnologie di rilevamento di presenza candidate. Le
prestazioni della colonna mmWave sono verificate sperimentalmente nel
\cref{cap:confronto}.}
\label{tab:comparativa}
\small
\begin{tabularx}{\textwidth}{@{}lXXX@{}}
\toprule
\textbf{Parametro} & \textbf{PIR} & \textbf{mmWave 24~GHz} & \textbf{UWB 3--10~GHz} \\
\midrule
Rilevamento corpo fermo & no (fisicamente) & sì (micro-movimenti) & sì \\
Misura di distanza & no & sì & sì \\
Conteggio persone & no & limitato (no angolazione) & limitato \\
Rilevamento respiro & no & sì, parziale & sì \\
Penetrazione ostacoli & no & parziale (non metallo/cemento) & buona \\
Corrente tipica & $\sim$0,05~mA & 50--80~mA & $<$10~mW \\
Privacy & assenza di dato personale & alta (nessuna immagine) & alta \\
Costo indicativo & 1--3~\euro & 5--15~\euro & 20--50~\euro \\
Alimentazione & batteria adeguata & fissa preferibile & batteria adeguata \\
\bottomrule
\end{tabularx}
\end{table}

Il quadro che emerge è quello di una \textbf{complementarità}, non di una
sostituzione: il PIR è imbattibile come guardiano a costo energetico nullo, il
mmWave è l'unico dei tre a offrire, a costo contenuto e con dati già elaborati,
una misura continua di \emph{quanto} si muove ciò che riflette. È esattamente la
grandezza necessaria all'indice di vitalità.

\medskip
\noindent
Come riferimento di prestazioni attese in condizioni ottimali, la letteratura
tecnica sui moduli mmWave consumer indica latenza di rilevamento del movimento
inferiore a 500~ms, conferma di persona ferma in 2--5~s, conferma di stanza vuota
in 20--60~s e tassi di falsi positivi inferiori a 0,05 eventi/h.
Il \cref{cap:confronto} verifica questi ordini di grandezza sul nostro esemplare.

\chapter{I moduli in esame e i dati che producono}
\label{cap:moduli}

\begin{quote}\small
Questo capitolo risponde all'obiettivo O2. Descrive i tre sensori effettivamente
usati, il protocollo con cui si interrogano, la piattaforma di acquisizione
realizzata per la campagna sperimentale e la natura dei dati prodotti, comprese le
loro insidie --- saturazione, canali strutturalmente disattivati, coda di presenza
--- individuate in fase di validazione.
\end{quote}

\section{HLK-LD2410B}
\label{sec:ld2410b}

\subsection{Specifiche}

Il \ldb\ è un modulo radar FMCW a 24~GHz prodotto da Shenzhen Hi-Link
Electronic. Le specifiche dichiarate dal manuale ufficiale e dal datasheet sono
riassunte nella \cref{tab:ld2410b-specs}.

\begin{table}[htbp]
\centering
\caption{Specifiche del \ldb\ (fonti: manuale ufficiale
V1.04~\cite{hilinkManualV104}, manuale V1.03~\cite{hilinkManualV103},
datasheet~\cite{ld2410bDatasheet}).}
\label{tab:ld2410b-specs}
\small
\begin{tabular}{@{}ll@{}}
\toprule
\textbf{Parametro} & \textbf{Valore} \\
\midrule
Tecnologia & FMCW, banda ISM 24~GHz \\
Portata dichiarata & 5--6~m \\
Campo visivo & $\pm 60\degree$ \\
Alimentazione & 5~V, alimentatore con capacità $>$200~mA \\
Corrente media & 80~mA \\
Interfaccia & UART 256000 baud, 8N1 \\
Bluetooth & integrato (variante B), password di default \texttt{HiLink} \\
Dimensioni & circa 35~mm $\times$ 7~mm \\
Gate di distanza & 9 (0--8) \\
Risoluzione gate & 0,75~m oppure 0,20~m \\
\bottomrule
\end{tabular}
\end{table}

\subsection{Piedinatura e cablaggio}
\label{sec:cablaggio}

Il modulo espone cinque piedini, definiti nella Tabella~1 del
datasheet~\cite{ld2410bDatasheet}:

\begin{center}\small
\begin{tabular}{@{}clll@{}}
\toprule
\textbf{Pin} & \textbf{Nome} & \textbf{Funzione} & \textbf{Collegamento ESP32} \\
\midrule
1 & OUT & presenza digitale (alto = persona) & non collegato \\
2 & UART\_Tx & il radar trasmette & GPIO25 (RX2) \\
3 & UART\_Rx & il radar riceve & GPIO26 (TX2) \\
4 & GND & massa & GND \\
5 & VCC & alimentazione & VIN (5~V) \\
\bottomrule
\end{tabular}
\end{center}

\begin{table}[htbp]
\centering
\caption{Corrispondenza fra i colori del cavo JST fornito con il modulo e i piedini
del datasheet. Una mappatura errata di questa corrispondenza ha comportato circa tre
settimane di lavoro perso nella convinzione che il modulo fosse guasto.}
\label{tab:colori}
\small
\begin{tabular}{@{}llll@{}}
\toprule
\textbf{Colore} & \textbf{Pin} & \textbf{Segnale} & \textbf{ESP32} \\
\midrule
Rosso & 5 & VCC & VIN (5~V) \\
Nero & 4 & GND & GND \\
Giallo & 3 & UART\_Rx & D26 (TX2) \\
Verde & 2 & UART\_Tx & D25 (RX2) \\
Blu & 1 & OUT & non collegato \\
\bottomrule
\end{tabular}
\end{table}

\paragraph{Nota metodologica.} Vale la pena riportare l'episodio, perché è un
risultato pratico di per sé. Il 19 luglio 2026 il modulo è stato dichiarato
guasto dopo giorni di tentativi in cui non si riceveva alcun byte; il 9 agosto,
ricontrollando la Tabella~1 del datasheet invece di dedurre l'ordine dei pin dalla
fotografia del connettore, si è scoperto che i due fili della seriale erano
invertiti. Il modulo ha risposto al primo tentativo. La regola che ne deriva ---
\emph{fare sempre riferimento alla tabella di piedinatura del datasheet, non
all'ordine apparente dei fili} --- è banale ma è costata tre settimane.

Un secondo caso analogo si è presentato con lo strumento di configurazione da PC:
il pacchetto \texttt{HLK-LD2420\_TOOL} non è utilizzabile con il LD2410B (protocolli
diversi, l'errore restituito è \emph{Failed to set data transfer mode}). Per il
LD2410B va usato il \texttt{LD2410 Tool} distribuito nella cartella Drive del
produttore~\cite{hilinkDriveLD2410B}.

\subsection{Protocollo seriale}
\label{sec:protocollo-seriale}

La comunicazione avviene su UART a 256000 baud, 8N1, con due tipi di frame
distinti da intestazione e chiusura (protocollo V1.07~\cite{hilinkProtoV107}, i
dettagli completi sono in \cref{app:protocollo}):

\begin{itemize}[leftmargin=1.4em,itemsep=2pt]
  \item \textbf{frame di comando} --- intestazione \texttt{FD FC FB FA}, chiusura
        \texttt{04 03 02 01};
  \item \textbf{frame di dati} --- intestazione \texttt{F4 F3 F2 F1}, chiusura
        \texttt{F8 F7 F6 F5}.
\end{itemize}

Nel frame di dati il byte 8 contiene lo stato del bersaglio come maschera di bit
(\texttt{0x01} = in movimento, \texttt{0x02} = stazionario), seguito dalle distanze
e dalle energie del bersaglio. Se il byte 6 vale \texttt{0x01} il frame è in
\emph{engineering mode} e trasporta anche le energie per-gate nei byte 19--37, più
la lettura del sensore di luce.

I comandi utilizzati in questo lavoro sono:

\begin{center}\small
\begin{tabular}{@{}ll@{}}
\toprule
\textbf{Comando} & \textbf{Funzione} \\
\midrule
\texttt{0x00FF} & entra in modalità configurazione \\
\texttt{0x00FE} & esce dalla modalità configurazione \\
\texttt{0x0062} & attiva engineering mode \\
\texttt{0x0063} & disattiva engineering mode \\
\texttt{0x00A0} & legge la versione firmware \\
\texttt{0x0061} & legge i parametri correnti (soglie, gate, timeout) \\
\texttt{0x0060} & imposta i parametri \\
\texttt{0x00AA} & imposta la risoluzione di distanza \\
\texttt{0x000B} & auto-calibrazione del rumore di fondo (firmware V2.44) \\
\bottomrule
\end{tabular}
\end{center}

\subsection{Engineering mode e dati per-gate}
\label{sec:engineering}

In modalità normale il modulo riporta, per ogni frame, lo stato di presenza, la
distanza e l'energia del bersaglio in movimento e di quello stazionario. In
\textbf{engineering mode} aggiunge le energie di ciascuno dei nove gate,
separatamente per la componente in movimento e per quella stazionaria: diciotto
valori aggiuntivi, più la lettura del fotosensore e lo stato del pin OUT.

Questo è il punto di svolta per gli obiettivi O2, O3 e O6: senza i dati per-gate
il radar sarebbe soltanto un PIR con la distanza, mentre con essi diventa uno
strumento di misura. Tutta l'analisi del respiro e l'indice di vitalità si basano
sulle serie temporali di questi canali.

\subsection{Identità dell'esemplare in prova}
\label{sec:identita}

Per la riproducibilità dei risultati si riportano i dati letti dal modulo
effettivamente usato, ottenuti via UART il 18 agosto 2026
(\cref{tab:identita}).

\begin{table}[htbp]
\centering
\caption{Identità e parametri di fabbrica dell'esemplare di \ldb\ in prova, letti
via UART.}
\label{tab:identita}
\small
\begin{tabular}{@{}ll@{}}
\toprule
\textbf{Dato} & \textbf{Valore letto} \\
\midrule
Versione firmware & 2.44.25070917 \\
Versione protocollo & 1 \\
MAC Bluetooth & 5E:E4:93:22:B9:3F \\
Baud confermato & 256000, 8N1 \\
Risoluzione gate & 75~cm; portata massima 675~cm (gate 0--8) \\
Timeout di presenza (\emph{no-one duration}) & 5~s \\
Soglie di movimento (gate 0--8) & 50, 50, 40, 30, 20, 15, 15, 15, 15 \\
Soglie stazionarie (gate 0--8) & 0, 0, 40, 40, 30, 30, 20, 20, 20 \\
\bottomrule
\end{tabular}
\end{table}

Il firmware 2.44 è l'ultima versione rilasciata dal produttore ed è quella che
introduce l'\textbf{auto-calibrazione del rumore di fondo}, cioè la taratura
automatica delle soglie di movimento e stazionarie~\cite{hilinkV244}.

\paragraph{Avvertenza sulla fonte dei parametri.} Lo strumento da PC
\texttt{LD2410 Tool} riporta per le sensibilità il valore~100 su tutti i gate. È
una lettura errata: la lettura via UART con il comando \texttt{0x0061} restituisce
la scala di fabbrica decrescente riportata in \cref{tab:identita}, che è coerente
sia con il protocollo V1.07 (\S1.2.2: il bersaglio è riconosciuto solo se
l'energia supera la soglia del gate) sia con il comportamento osservato. Per i
valori numerici si è quindi fatto fede alla lettura via UART e non allo strumento
da PC.

\section{HLK-LD2420}
\label{sec:ld2420}

Il \ldc\ è il secondo radar 24~GHz in dotazione; la documentazione di
riferimento è quella distribuita dal produttore~\cite{hilinkDriveLD2420}. Le differenze rilevanti rispetto
al LD2410B sono riassunte nella \cref{tab:confronto-moduli}. Due meritano
attenzione particolare:

\begin{itemize}[leftmargin=1.4em,itemsep=2pt]
  \item \textbf{alimentazione a 3,3~V} e non a 5~V: collegarlo come il LD2410B lo
        danneggerebbe;
  \item \textbf{la piedinatura cambia con la versione firmware}. Fino alla 1.5.2 il
        piedino 3 (OT1) è l'uscita di presenza digitale e il piedino 5 (OT2) è il
        TX seriale; dalla 1.5.3 i ruoli sono invertiti, e cambia anche il baud
        (256000 fino alla 1.5.2, 115200 dalla 1.5.3). Verificare la versione prima
        di collegare non è un consiglio ma un prerequisito.
\end{itemize}

\begin{table}[htbp]
\centering
\caption{Confronto fra i due moduli radar in dotazione.}
\label{tab:confronto-moduli}
\small
\begin{tabularx}{\textwidth}{@{}lXX@{}}
\toprule
\textbf{Caratteristica} & \textbf{LD2410B} & \textbf{LD2420} \\
\midrule
Alimentazione & 5~V & 3,3~V \\
Portata & circa 5--6~m & circa 12~m \\
Gate & 9 (0--8) & 15 (0--14) \\
Risoluzione gate & 0,75~m oppure 0,20~m & 0,70~m fisso \\
Baud & 256000 fisso & 115200 o 256000 (dipende dal firmware) \\
Bluetooth & presente & assente \\
Calibrazione & manuale (auto da V2.44) & automatica (firmware $\geq$ 1.5.4) \\
Dati per-gate & molto dettagliati & limitati \\
Piedinatura & fissa & varia con il firmware \\
Uso più adatto & respiro, dettaglio & portata lunga, copertura \\
\bottomrule
\end{tabularx}
\end{table}

\subsection{Verifica sull'esemplare in dotazione}

L'esemplare in prova è stato collegato all'ESP32 (GPIO16/17) e verificato il
17 luglio 2026. Risultati:

\begin{itemize}[leftmargin=1.4em,itemsep=2pt]
  \item il baud funzionante è \textbf{115200}, da cui si deduce firmware
        $\geq$ 1.5.3 e quindi TX seriale sul piedino 3 (OT1). A 256000 baud si
        riceve solo rumore, con lo stesso alfabeto di byte tipico di un log letto
        al baud sbagliato;
  \item la modalità di fabbrica è \textbf{ASCII ``semplice''} e non il protocollo
        binario: il modulo trasmette righe di testo terminate da
        \texttt{CR\,LF} contenenti \texttt{ON}, \texttt{OFF} oppure
        \texttt{Range NN};
  \item in questa modalità sono disponibili solo presenza e un valore di
        ``range'': \textbf{nessuna energia per-gate}, nessuna distinzione fra
        movimento e stazionario.
\end{itemize}

\paragraph{Limite dichiarato.} L'unità di misura del campo \texttt{Range} non è
documentata nella documentazione in nostro possesso. I valori osservati
(7--37 muovendosi nella stanza) sono compatibili con decimetri, cioè 0,7--3,7~m, ma
in assenza di conferma il dato non viene usato come misura di distanza in questa
tesi. Per accedere all'engineering mode del LD2420 occorre riconfigurare il modulo
in modalità binaria; l'attività è rimasta fuori dal perimetro sperimentale
completato (\cref{sec:ld2420-risultati}).

\section{Il sensore PIR in dotazione: HC-SR501}
\label{sec:hcsr501}

Il modulo in dotazione è stato identificato per via fotografica come un
\textbf{\pir}, riconoscendo sul circuito stampato la serigrafia del circuito di
condizionamento BISS0001 e del regolatore lineare HT7133. Le specifiche da
datasheet~\cite{hcsr501} sono riassunte nella \cref{tab:pir-specs}.

\begin{table}[htbp]
\centering
\caption{Specifiche del PIR \pir\ in dotazione (datasheet~\cite{hcsr501}).}
\label{tab:pir-specs}
\small
\begin{tabularx}{\textwidth}{@{}lX@{}}
\toprule
\textbf{Parametro} & \textbf{Valore} \\
\midrule
Base & BISS0001 + elemento LHI778 \\
Alimentazione & 4,5--20~V DC (collegato a VIN 5~V dell'ESP32) \\
Livello logico di uscita & alto 3,3~V / basso 0~V, grazie al regolatore HT7133 a
  bordo: compatibile con i GPIO dell'ESP32 senza adattamento \\
Corrente di riposo & $<$50~\textmu A \\
Portata & 3--7~m, regolabile con il trimmer di sensibilità \\
Angolo di copertura & cono $<$110\degree\ (più ampio dei $\pm 60\degree$ del radar) \\
Tempo di ritenuta & regolabile da circa 3~s a circa 5~min \\
Block time & 2,5~s dopo ogni rilascio \\
Trigger & jumper: L = singolo, H = ripetibile \\
Temperatura operativa & $-15$ / $+70$~\celsius \\
Dimensioni & 32 $\times$ 24~mm \\
\bottomrule
\end{tabularx}
\end{table}

\paragraph{Configurazione adottata per la campagna} (fissata il 19 luglio 2026 e
mantenuta identica per tutte le acquisizioni, con fotografia della posizione dei
trimmer):

\begin{enumerate}[leftmargin=1.8em,itemsep=2pt]
  \item jumper in posizione \textbf{H} (repeat trigger), altrimenti una persona in
        movimento continuo apparirebbe intermittente per costruzione;
  \item trimmer del tempo di ritenuta al \textbf{minimo}: la ritenuta misurata è
        risultata di circa 3,4~s;
  \item trimmer di sensibilità a \textbf{metà corsa};
  \item uscita collegata a \textbf{GPIO34}. La scelta merita una nota: GPIO34 è un
        piedino di solo ingresso e privo di resistenza di pull-up interna, il che
        normalmente lo rende inadatto a un sensore a collettore aperto;
        l'\pir{} però pilota attivamente la propria uscita a 3,3~V e 0~V, quindi il
        collegamento è corretto. GPIO25 non era disponibile perché usato come RX2
        del radar.
\end{enumerate}

\section{La piattaforma di acquisizione}
\label{sec:piattaforma}

\subsection{Punto di partenza: il repository di riferimento}

Il lavoro parte dal repository del relatore~\cite{repoCallisto}, che contiene un
firmware PlatformIO per ESP32 basato sulla libreria \texttt{ncmreynolds/ld2410} e
uno script Python di acquisizione. Il firmware legge il radar via UART, gestisce
opzionalmente un display OLED e un PIR e stampa righe CSV sulla seriale; lo script
\texttt{acquire.py} le salva su file aggiungendo i metadati sperimentali passati
da riga di comando.

Lo studio del codice ha evidenziato tre limiti per gli scopi di questa tesi:

\begin{enumerate}[leftmargin=1.8em,itemsep=3pt]
  \item \textbf{campionamento a 1~Hz}. Troppo lento per misurare latenze con
        precisione e, soprattutto, insufficiente per l'analisi del respiro: con
        frequenza di campionamento 1~Hz il limite di Nyquist è 0,5~Hz, cioè
        esattamente il bordo superiore della banda respiratoria;
  \item \textbf{assenza dell'engineering mode}: viene registrata solo l'energia del
        bersaglio, non i diciotto valori per-gate;
  \item \textbf{piedinatura opposta} alla nostra (radar TX su GPIO17 e RX su
        GPIO16), quindi i due firmware non sono interscambiabili senza riscablare.
\end{enumerate}

\subsection{Il firmware realizzato}

È stato quindi scritto un logger dedicato
(\texttt{firmware/ld2410b\_logger}, \cref{app:codice}) basato sulla libreria
\texttt{MyLD2410}, con le seguenti caratteristiche:

\begin{itemize}[leftmargin=1.4em,itemsep=2pt]
  \item campionamento a \textbf{5~Hz} con periodo fisso;
  \item \textbf{engineering mode} sempre attivo, con registrazione di tutti i
        canali;
  \item lettura del PIR su GPIO34 nello stesso frame, in modo che il confronto fra i
        due sensori sia \emph{sincrono per costruzione} --- non due acquisizioni da
        allineare a posteriori;
  \item formato CSV retrocompatibile con \texttt{acquire.py}.
\end{itemize}

Il formato CSV, unico per tutta la tesi, è il seguente (le prime nove colonne
coincidono con quelle del repository di riferimento):

\begin{lstlisting}[caption={Intestazione del CSV prodotto dal logger, in
engineering mode.},label={lst:csv},basicstyle=\ttfamily\scriptsize]
timestamp_ms, radar_presence, moving_target, stationary_target,
moving_distance_cm, stationary_distance_cm, moving_energy, stationary_energy,
pir_presence,
menergy_gate0 ... menergy_gate8,   (9 colonne)
senergy_gate0 ... senergy_gate8,   (9 colonne)
light_level, out_level
\end{lstlisting}

Lo script di acquisizione aggiunge a ogni riga i metadati sperimentali:
\texttt{pc\_time\_s}, \texttt{group\_id}, \texttt{trial\_id}, \texttt{scenario},
\texttt{ground\_truth\_presence}, \texttt{ground\_truth\_state}.

\paragraph{Una correzione necessaria allo script di acquisizione.} La versione
originale scrive righe sul file solo dopo aver riconosciuto la riga di intestazione,
che l'ESP32 stampa una sola volta in \texttt{setup()}. Se il reset all'apertura
della porta seriale non scatta --- cosa che accade con una certa frequenza --- il
CSV risulta \textbf{vuoto senza alcun messaggio di errore}, e lo si scopre a
sessione finita. La nostra versione ricostruisce l'intestazione dal numero di
colonne della prima riga di dati (9 = base, 27 = engineering, 29 = engineering con
luce e OUT). È una modifica di poche righe che ha evitato la perdita di sessioni
lunghe.

\subsection{La catena di analisi}
\label{sec:catena-analisi}

L'analisi è svolta da tre script Python, scritti per questa tesi e basati sulla
sola libreria standard, tranne dove indicato:

\begin{description}[leftmargin=1.4em,style=nextline,itemsep=3pt]
  \item[\texttt{analisi/verifica\_engineering.py}] controllo di qualità di un CSV
    \emph{prima} di usarlo: cadenza reale e jitter, presenza delle colonne per-gate,
    percentuale di saturazione a 100, coerenza fra gate di picco e distanza
    riportata, rumore di fondo per-gate a stanza vuota, transizioni di presenza. Va
    eseguito a ogni sessione: è ciò che ha permesso di individuare i problemi
    descritti in \cref{sec:insidie}.
  \item[\texttt{analisi/analizza\_test.py}] metriche di confronto: accuratezza,
    falsi positivi in eventi/h, falsi negativi in percentuale, latenza di
    rilevamento e di rilascio, statistiche di distanza ed energia, aggregazione per
    scenario con media e deviazione standard sui cinque trial.
  \item[\texttt{analisi/analizza\_respiro.py}] analisi in frequenza: FFT della
    serie di energia, individuazione del picco in banda 0,1--0,5~Hz, stima degli
    atti al minuto, esportazione dello spettro in CSV. Richiede NumPy. La modalità
    \texttt{-{}-scan} prova tutti e venti i canali di energia, scarta quelli saturi
    e quelli piatti, li ordina per rapporto segnale-rumore e riporta la mediana
    delle stime concordi.
\end{description}

L'ultima funzione non era prevista inizialmente: è nata dal fatto che nel primo
tentativo di misura del respiro il canale scelto per intuito
(\texttt{stationary\_energy}) era completamente saturo, mentre il segnale era
perfettamente leggibile su \texttt{menergy\_gate2}. La ricerca automatica del
canale migliore, con criterio dichiarato, è diventata parte del metodo.

\section{Che aspetto hanno i dati: validazione e insidie}
\label{sec:insidie}

Prima di usare i dati per il confronto quantitativo è stata svolta una sessione di
validazione (673 campioni su 134~s). I risultati positivi sono:

\begin{itemize}[leftmargin=1.4em,itemsep=2pt]
  \item \textbf{cadenza di 5,00~Hz esatti}: intervallo di 200~ms su tutti i 672
        intervalli, jitter nullo. La serie temporale è uniforme, requisito
        necessario per l'analisi in frequenza;
  \item \textbf{engineering mode attivo}, con tutte e diciotto le colonne per-gate
        popolate;
  \item \textbf{mappa gate--distanza validata al 97,1\,\%}: la percentuale di
        campioni in cui il gate di massima energia coincide, entro $\pm 1$, con
        quello atteso dalla distanza riportata (distanza divisa per 0,75~m), su
        tutto il range 0--6~m. Su una seconda sessione indipendente la conferma è
        salita al 98,3\,\% su 663 campioni;
  \item \textbf{zero falsi positivi} nei 29~s di stanza vuota della sessione.
\end{itemize}

La stessa validazione ha però messo in luce tre comportamenti del modulo che
condizionano tutte le analisi successive e che vanno dichiarati esplicitamente.

\subsection{Saturazione dei canali di energia}
\label{sec:saturazione}

Con il bersaglio fermo a distanza ravvicinata, l'energia stazionaria è a fondo
scala (valore 100) nel \textbf{92\,\%} dei campioni e quella in movimento nel
\textbf{41\,\%}. Un segnale che tocca il fondo scala non porta informazione:
qualunque analisi di variazione su un canale saturo è priva di senso.

Va chiarito un punto che è stato oggetto di un'ipotesi errata iniziale: \textbf{le
soglie non risolvono la saturazione}. L'energia è un valore normalizzato 0--100 e
la soglia interviene soltanto sulla \emph{decisione} di rilevamento (protocollo
V1.07, \S1.2.2), non sul valore riportato. Nemmeno l'auto-calibrazione del firmware
V2.44 la risolve, perché anch'essa tara soglie e non la scala. L'unico rimedio è
\textbf{geometrico}: allontanare il soggetto (almeno 1,5--2~m), disallineare il
sensore o attenuare il segnale. Questo vincolo ha determinato il progetto degli
esperimenti sul respiro e sulla vitalità.

\subsection{Canali stazionari strutturalmente disattivati sotto 1,5~m}
\label{sec:senergy-zero}

Le colonne \texttt{senergy\_gate0} e \texttt{senergy\_gate1} sono risultate pari a
zero in tutti i 673 campioni. Non è un difetto dell'esemplare: anche l'esempio
riportato nel protocollo ufficiale V1.07 (\S2.3.2) mostra \texttt{00 00} per i
primi due gate stazionari, e le relative soglie di fabbrica valgono~0
(\cref{tab:identita}) --- con la regola ``energia \emph{maggiore} della soglia'' il
canale è quindi disattivato per costruzione sotto 1,5~m.

È importante non trarre la conclusione sbagliata. Il radar \textbf{riporta comunque}
bersagli fermi sotto 1,5~m: nei dati si osservano valori di
\texttt{stationary\_distance} fra 70 e 89~cm con \texttt{stationary\_energy} pari a
100. Ciò che manca è solo il \emph{dettaglio per-gate} dello stazionario nei primi
due gate. La conseguenza pratica è però rilevante: le analisi che si basano sulle
serie per-gate --- respiro e indice di vitalità --- per un soggetto a meno di 1,5~m
devono usare i \textbf{canali in movimento} oppure i gate di indice $\geq 2$.
E poiché nello scenario UPRISE la persona è sotto un banco, cioè vicina, questo è
un vincolo di progetto e non un dettaglio (\cref{sec:sotto-banco}).

\subsection{Coda di presenza}
\label{sec:coda}

Dopo l'uscita del soggetto dalla stanza il radar ha mantenuto
\texttt{radar\_presence = 1} per circa 10~s, riportando un bersaglio fermo
``fantasma'' a circa 5,2~m con energia in decadimento. Il comportamento è atteso ed
è il \emph{no-one duration} previsto dal protocollo (\S1.2.2), impostato a 5~s di
fabbrica sul nostro esemplare.

Questo ha una conseguenza metodologica precisa: la coda va misurata come
\textbf{latenza di rilascio} (\cref{sec:latenza}) e non va conteggiata come falso
positivo. Confondere le due cose porterebbe ad attribuire al radar decine di falsi
positivi inesistenti.

\subsection{Il pin OUT e il sensore di luce}

Due verifiche minori ma utili. La colonna \texttt{out\_level}, che riporta lo stato
del pin OUT letto \emph{dal frame UART}, coincide con \texttt{radar\_presence} in
1699 campioni su 1699: il collegamento fisico del filo blu è quindi inutile ai
nostri fini. È però un'informazione utile per il progetto UPRISE, perché indica che
un dispositivo interessato alla sola presenza binaria può usare il pin OUT senza
implementare la UART.

Il sensore di luce funziona: la colonna \texttt{light\_level} vale 21--29 di giorno
e 0--1 di notte, misurato su una sessione notturna di 6,5~h. Il dato è annotato
perché il comando di configurazione ausiliaria \texttt{0x01AE} non risponde sul
nostro esemplare: i due fatti sono indipendenti, il valore fotosensibile viaggia
nel frame di engineering mode e segue effettivamente l'illuminazione ambientale.

\section{Decisione sulle soglie: non ricalibrare}
\label{sec:soglie}

Un'ultima scelta metodologica va motivata, perché condiziona la validità di tutti i
confronti fra trial. Il \cref{tab:soglie} confronta le soglie di fabbrica con il
rumore di fondo misurato a stanza vuota.

\begin{table}[htbp]
\centering
\caption{Soglie di fabbrica per-gate confrontate con il rumore di fondo misurato a
stanza vuota (29~s). Regola del protocollo V1.07 \S1.2.2: il bersaglio è
riconosciuto solo se l'energia supera la soglia.}
\label{tab:soglie}
\small
\begin{tabular}{@{}clrrrrr@{}}
\toprule
\textbf{Gate} & \textbf{Distanza} & \multicolumn{3}{c}{\textbf{Movimento}} &
\multicolumn{2}{c}{\textbf{Stazionario}} \\
\cmidrule(lr){3-5}\cmidrule(lr){6-7}
 & & soglia & rumore max & margine & soglia & rumore max \\
\midrule
0 & 0--75~cm    & 50 & 26 & 24 & 0  & 0 \\
1 & 75--150~cm  & 50 & 21 & 29 & 0  & 0 \\
2 & 150--225~cm & 40 & 9  & 31 & 40 & 6 \\
3 & 225--300~cm & 30 & 4  & 26 & 40 & 4 \\
4 & 300--375~cm & 20 & 9  & 11 & 30 & 4 \\
5 & 375--450~cm & 15 & 6  & 9  & 30 & 3 \\
6 & 450--525~cm & 15 & 11 & \textbf{4} & 20 & 5 \\
7 & 525--600~cm & 15 & 8  & 7  & 20 & 7 \\
8 & 600--675~cm & 15 & 7  & 8  & 20 & 5 \\
\bottomrule
\end{tabular}
\end{table}

Il margine più stretto è il gate~6 in movimento (soglia 15 contro rumore massimo
11), che è anche il gate in cui si è manifestato il bersaglio fantasma della coda di
presenza. Su tutti gli altri gate il margine è ampio, e i falsi positivi osservati
sono zero. Si è quindi deciso di \textbf{non ricalibrare}: la taratura di fabbrica
regge nell'ambiente di prova, e le soglie restano il \emph{riferimento fisso} per
tutta la campagna. Cambiarle a metà campagna renderebbe i trial non confrontabili
fra loro, che è il difetto peggiore in un lavoro comparativo.

L'auto-calibrazione disponibile nel firmware V2.44 va usata dopo, come esperimento a
sé stante con confronto prima/dopo: è materiale per una sezione della tesi, non un
prerequisito della campagna.

La stabilità di questa scelta è stata verificata su tempi lunghi: nella sessione
notturna di 6,5~h la media dell'energia di rumore del gate~0 è passata da 17,4 a
17,8 fra la prima e la seconda metà della notte e quella del gate~1 da 13,3 a 13,1,
con un massimo assoluto di 34 sul gate~0 contro una soglia di 50.

\chapter{Confronto sperimentale PIR vs mmWave}
\label{cap:confronto}

\begin{quote}\small
Questo è il capitolo centrale della tesi e risponde all'obiettivo O3. Presenta il
protocollo di misura, i risultati acquisiti e la loro discussione critica. Le
sezioni contrassegnate come \emph{in corso} indicano acquisizioni pianificate ma
non ancora completate alla data di questa stesura: la loro collocazione nel
capitolo è già definita, così come le metriche e i comandi con cui verranno
prodotte.
\end{quote}

\section{Protocollo sperimentale}
\label{sec:protocollo}

\subsection{Principi generali}

Il protocollo è stato definito \emph{prima} di iniziare le acquisizioni ed è
rimasto invariato, per una ragione precisa: la campagna completa richiede circa
dieci ore di acquisizioni, e una modifica del protocollo a metà strada le renderebbe
non confrontabili.

I principi adottati sono cinque.

\begin{enumerate}[leftmargin=1.8em,itemsep=4pt]
  \item \textbf{Cinque ripetizioni per scenario} (T01--T05). È ciò che permette di
        riportare media e deviazione standard, e quindi di rendere il confronto
        ``numerico'' e non anecdottico. Ogni valore riportato in questo capitolo
        nella forma $\mu \pm \sigma$ è calcolato su cinque trial indipendenti.
  \item \textbf{Acquisizione sincrona dei due sensori}. Radar e PIR sono letti nello
        stesso frame dal medesimo firmware, sulla stessa base dei tempi: il
        confronto non richiede allineamento a posteriori e non introduce errori di
        sincronizzazione.
  \item \textbf{Ground truth dichiarata}. Ogni acquisizione porta nei metadati la
        presenza reale e lo stato reale del soggetto, dichiarati dall'operatore da
        riga di comando. Per la ground truth statica per file questo è sufficiente;
        per le latenze si adotta la convenzione dell'evento a istante noto
        (\cref{sec:latenza}).
  \item \textbf{Setup fisso}. LD2410B e PIR affiancati sul bordo di un tavolo a
        circa 1~m di altezza, puntati verso l'area di prova, con la zona di prova
        interna al campo di \emph{entrambi} i sensori (il PIR copre un cono più
        ampio, $<$110\degree, del radar). Nessun oggetto in movimento nella stanza:
        nessun ventilatore, tende ferme, porta chiusa.
  \item \textbf{Alimentazione da alimentatore di rete} da almeno 1~A e non dalla
        porta USB del portatile. L'ESP32 con WiFi attivo e il LD2410B presentano
        picchi superiori a 400~mA, e una porta USB debole provoca brownout, cioè
        riavvii casuali che si presentano come bug del firmware.
\end{enumerate}

\subsection{Metriche}

Le metriche calcolate da \texttt{analisi/analizza\_test.py} sono le seguenti.

\begin{description}[leftmargin=1.4em,style=nextline,itemsep=3pt]
  \item[Falsi negativi (\%)] percentuale dei campioni con presenza reale in cui il
    sensore riporta assenza. È la metrica chiave dello scenario ``persona ferma''.
  \item[Falsi positivi (eventi/h)] numero di transizioni $0 \to 1$ per ora di
    osservazione a stanza vuota. Si conta per \emph{eventi} e non per campioni,
    perché un singolo falso positivo che dura dieci secondi non è dieci volte più
    grave di uno che dura un secondo.
  \item[Latenza di rilevamento] tempo fra l'evento reale e la prima transizione
    $0 \to 1$ del sensore.
  \item[Latenza di rilascio] tempo fra l'uscita del soggetto e la transizione
    stabile $1 \to 0$, con conteggio separato delle eventuali riaccensioni nella
    coda.
  \item[Accuratezza della distanza] scostamento fra distanza riportata dal radar e
    distanza reale misurata, con separazione fra dispersione \emph{fra} trial e
    dispersione \emph{entro} il trial.
\end{description}

\subsection{Registro delle sessioni}
\label{sec:registro}

Ogni sessione è annotata in un registro con data, ora, temperatura della stanza,
soggetto, alimentazione, test svolti e note. La temperatura è registrata perché il
PIR degrada al crescere della temperatura ambiente (\cref{sec:pir-temperatura}); le
sessioni di questa campagna si collocano fra 25 e 29\,\celsius, condizione
sfavorevole al PIR che va tenuta presente nella discussione (\cref{sec:limiti}).

Le acquisizioni completate sono riassunte nella \cref{tab:sessioni}.

\begin{table}[htbp]
\centering
\caption{Sessioni di acquisizione completate.}
\label{tab:sessioni}
\small
\begin{tabularx}{\textwidth}{@{}llrX@{}}
\toprule
\textbf{Data} & \textbf{Scenario} & \textbf{Durata} & \textbf{Note} \\
\midrule
18/08/2026 & validazione logger & 134~s & 673 campioni, 5,00~Hz esatti \\
18/08/2026 & pilota PIR vs radar & 339,6~s & 1699 campioni, PIR su GPIO34 \\
18/08/2026 & pilota respiro & 91~s & soggetto a 40~cm, nessuna ground truth \\
18/08/2026 & stanza vuota (giorno) & 1742~s & 8711 campioni, 29\,\celsius \\
18/08/2026 & fermo seduto T01--T05 & 5 $\times$ 302~s & 7554 campioni, 2,30~m, 27\,\celsius \\
20/08/2026 & distanza 1--5~m, 25 trial & 25 $\times$ 62~s & cammino sul posto \\
20--21/08/2026 & stanza vuota (notte) & 23640,6~s & 118\,204 campioni, 6,57~h \\
\bottomrule
\end{tabularx}
\end{table}

\paragraph{Nota su una sessione troncata.} La sessione notturna era programmata per
sette ore ed è terminata a 6,57~h perché un aggiornamento automatico di Windows ha
riavviato il computer. I dati restano validi: l'intervallo di campionamento è di
200~ms esatti su tutti i 118\,203 intervalli, senza buchi né campioni persi ---
l'acquisizione è semplicemente troncata. L'episodio è però istruttivo e ha prodotto
una regola di protocollo: per le acquisizioni oltre l'ora, sospendere gli
aggiornamenti di Windows e verificare che sospensione e ibernazione siano disattivate.

\section{Falsi positivi: stanza vuota}
\label{sec:falsi-positivi}

\paragraph{Procedura.} Acquisizione avviata con l'operatore che esce immediatamente
dalla stanza; scarto dei primi 60~s (stabilizzazione del PIR e uscita
dell'operatore); conteggio delle transizioni $0 \to 1$ sul resto.

\paragraph{Risultato.} Su un totale di \textbf{7,05~h} di stanza vuota
(6,57~h notturni più 0,48~h diurni), \textbf{né il radar né il PIR hanno prodotto
alcun falso positivo}. Nella sessione notturna gli unici 96 campioni con presenza
sono compresi fra $t=0$ e $t=19$~s e corrispondono all'operatore che esce dalla
stanza.

Con zero eventi osservati non è corretto scrivere ``0 eventi/h'': si riporta il
limite superiore, che dipende dal tempo di osservazione. Applicando la regola del
tre ($3/T$, che approssima il limite superiore al 95\,\% di confidenza per un
processo di Poisson con zero eventi osservati):

\begin{equation}
\lambda_{95\%} \leq \frac{3}{7{,}05~\text{h}} = 0{,}43~\text{eventi/h}
\label{eq:regola-tre}
\end{equation}

Il risultato da riportare è quindi: \textbf{falsi positivi $\leq 0{,}43$ eventi/h
per entrambi i sensori, al 95\,\% di confidenza}. Il valore è coerente con l'ordine
di grandezza indicato dalla letteratura tecnica per i moduli mmWave consumer
($<$0,05 eventi/h), che però richiederebbe una sessione di decine di ore per essere
confermato con questa metodologia. Questo è il modo corretto di riportare
un'assenza di eventi, e la principale limitazione del risultato è il tempo di
osservazione, non il sensore.

\paragraph{Nota sulla temperatura.} Le sessioni sono state svolte a 29\,\celsius. È
una condizione \emph{sfavorevole} al PIR per il rilevamento (minore contrasto
termico) ma \emph{favorevole} sui falsi positivi: un ambiente caldo e uniforme
produce meno gradienti spuri. Il risultato ``zero falsi positivi PIR'' va quindi
letto come non pessimistico.

\section{Il risultato centrale: la persona ferma}
\label{sec:persona-ferma}

\paragraph{Procedura.} Soggetto seduto immobile a 2,30~m dalla coppia di sensori,
respirazione spontanea, nessun movimento volontario. Cinque trial da 302~s utili
ciascuno dopo lo scarto dei primi 30~s. Temperatura della stanza 27\,\celsius.
Totale: 7554 campioni con presenza reale.

\paragraph{Risultato.} I valori misurati sono riportati nella
\cref{tab:fermo-seduto}.

\begin{table}[htbp]
\centering
\caption{Scenario ``persona ferma seduta a 2,30~m''. Falsi negativi sui 7554
campioni con presenza reale, su cinque trial indipendenti da 302~s.}
\label{tab:fermo-seduto}
\small
\begin{tabular}{@{}lrr@{}}
\toprule
\textbf{Metrica} & \textbf{PIR \pir} & \textbf{Radar \ldb} \\
\midrule
Falsi negativi (media $\pm$ dev.\ std.\ su 5 trial) & $99{,}78 \pm 0{,}49\,\%$ &
  $\mathbf{0{,}00 \pm 0{,}00\,\%}$ \\
Trial con falsi negativi al 100\,\% & 4 su 5 & 0 su 5 \\
Campioni con presenza rilevata & $\sim$17 su 7554 & 7554 su 7554 \\
\bottomrule
\end{tabular}
\end{table}

Il dato va letto con attenzione al suo significato operativo: in quattro trial su
cinque il PIR ha prodotto \textbf{zero rilevamenti in cinque minuti} con una
persona viva seduta a 2,3~m di distanza, in campo visivo libero. Il radar ha
rilevato la presenza nel 100\,\% dei campioni in tutti e cinque i trial.

Questo è il risultato che risponde alla domanda di ricerca principale. Non è un
esito sorprendente --- è precisamente ciò che la fisica del piroelettrico prevede
(\cref{sec:pir-principio}) --- ma la tesi lo \emph{quantifica}, con deviazione
standard su cinque ripetizioni, nella condizione che interessa il progetto UPRISE.

Un risultato analogo, ottenuto in una sessione precedente con procedura meno
controllata, conferma il quadro: in un tratto di \textbf{93,2~s di immobilità
continuativa} a circa 2,9~m, il radar ha rilevato la presenza nel 100\,\% dei
campioni e il PIR nello 0,2\,\% (un solo campione su 466).

\subsection{Come va interpretato il valore percentuale}
\label{sec:interpretazione-pir}

Il numero ``99,78\,\% di falsi negativi'' richiede una precisazione onesta, perché
dipende in parte da un parametro sotto il nostro controllo: il trimmer del tempo di
ritenuta, che abbiamo tenuto al minimo (\cref{sec:hcsr501}).

Nella sessione pilota il PIR ha prodotto \textbf{14 impulsi di durata costante
compresa fra 3,4 e 3,6~s} invece di un segnale continuo. Su due sessioni sono stati
misurati in totale \textbf{30 impulsi, tutti compresi fra 3,4 e 3,8~s}. Questa è la
prova quantitativa di un fatto qualitativo: \textbf{l'uscita del PIR è un
monostabile a durata fissa}, indipendente da cosa fa la persona. Il PIR non misura
la presenza né la durata del movimento: segnala eventi.

La conseguenza è che, con la ritenuta regolata al massimo, il PIR
\emph{sembrerebbe} rilevare molto di più --- ma starebbe soltanto trattenendo
l'ultimo evento per minuti. La percentuale di campioni rilevati non è quindi una
proprietà robusta del sensore.

Il dato robusto e non contestabile è l'altro, e va riportato come tale: nei 93~s di
immobilità continuativa e nei cinque trial da cinque minuti, \textbf{il PIR non ha
prodotto alcun evento di movimento} (con le eccezioni contate sopra), perché non
c'era alcun transito di flusso infrarosso da rilevare. La percentuale di campioni è
il modo di quantificarlo, il numero di eventi è il modo di dimostrarlo.

\section{Il PIR sbaglia anche con la persona in movimento}
\label{sec:pir-movimento}

Questo è il risultato meno atteso della campagna, e con il maggiore impatto sul
contesto applicativo.

\paragraph{Procedura.} Soggetto che cammina \emph{sul posto} a distanza fissa dalla
coppia di sensori --- movimento continuo, ma senza spostamento nello spazio. Cinque
distanze (1, 2, 3, 4, 5~m), cinque trial per distanza, 62~s utili ciascuno: 25
trial in totale.

\paragraph{Risultato.} A 1~m, condizione teoricamente ideale per un PIR (bersaglio
vicino, movimento continuo, campo libero), il PIR ha registrato
$\mathbf{80{,}46 \pm 6{,}22\,\%}$ di falsi negativi. In 62~s di cammino continuo ha
emesso da \textbf{1 a 4 impulsi} per trial, di durata 3,6--3,8~s. A 2, 3, 4 e 5~m il
tasso di falsi negativi è \textbf{100,0\,\%} su tutti e venti i trial: nessun
rilevamento. Il radar ha rilevato la presenza nel 100\,\% dei campioni a tutte e
cinque le distanze.

\paragraph{Spiegazione.} Il risultato non è un'anomalia e si spiega esattamente con
il meccanismo della lente di Fresnel (\cref{sec:fresnel}): il piroelettrico
risponde al \emph{transito} del flusso infrarosso fra le zone della lente, non al
movimento in sé. Chi si muove \emph{sul posto} non attraversa le zone, quindi non
genera transizioni --- e per il sensore è indistinguibile da un corpo fermo.

\paragraph{Rilevanza per UPRISE.} Questo risultato è più importante di quanto la
sua sobrietà suggerisca. Una persona intrappolata sotto un arredo \emph{si muove sul
posto}: si agita, cerca di liberarsi, sposta le braccia, chiama aiuto. Non attraversa
la stanza. È il caso peggiore per un PIR, e la campagna mostra che il PIR fallisce
in questo caso quasi allo stesso modo in cui fallisce sulla persona completamente
immobile ($80\,\%$ contro $\sim\!100\,\%$ di falsi negativi, e $100\,\%$ oltre i 2~m).

Detto in altri termini: l'ipotesi implicita ``il PIR non vede chi è fermo ma vede
chi si muove'' è, nello scenario di interesse, falsa.

\section{Accuratezza della misura di distanza}
\label{sec:distanza}

Gli stessi 25 trial forniscono la caratterizzazione della misura di distanza del
radar, riportata nella \cref{tab:distanza}.

\begin{table}[htbp]
\centering
\caption{Accuratezza della distanza riportata dal \ldb, su 25 trial (5 distanze
$\times$ 5 ripetizioni). L'energia è il valore medio del bersaglio, normalizzato
0--100.}
\label{tab:distanza}
\small
\begin{tabular}{@{}rrrrr@{}}
\toprule
\textbf{Distanza reale} & \textbf{Errore grezzo} & \textbf{Errore relativo} &
\textbf{Dispersione entro trial} & \textbf{Energia media} \\
\midrule
1~m & $-0{,}5$~cm & $-0{,}5\,\%$ & 10--26~cm & 99 \\
2~m & --- & --- & 10--26~cm & 85 \\
3~m & --- & --- & 10--26~cm & 54 \\
4~m & --- & --- & 10--26~cm & 35 \\
5~m & $+18{,}2$~cm & $+3{,}6\,\%$ & 10--26~cm & 28 \\
\bottomrule
\end{tabular}
\end{table}

\todo{completare le colonne intermedie della \cref{tab:distanza} con i valori
puntuali per 2, 3 e 4~m e la dispersione entro trial per singola distanza,
rileggendoli dall'output di \texttt{analizza\_test.py}}

\subsection{Il modello di errore}

La regressione lineare fra distanza misurata e distanza reale, sui 25 trial, dà:

\begin{equation}
d_{\text{misurata}} = 1{,}0381 \cdot d_{\text{reale}} - 1{,}32~\text{cm},
\qquad R^2 = 0{,}99965
\label{eq:regressione}
\end{equation}

Il risultato si legge in due parti, ed è la lettura --- non il numero grezzo --- il
contributo di questa sezione.

\begin{itemize}[leftmargin=1.4em,itemsep=3pt]
  \item Il radar è \textbf{estremamente lineare}: $R^2 = 0{,}99965$ e residui dalla
        retta non superiori a 4,9~cm su tutto il range. Non c'è alcuna non linearità
        apprezzabile fra 1 e 5~m.
  \item L'errore è quasi interamente un \textbf{errore di scala del $+3{,}81\,\%$},
        con offset praticamente nullo ($-1{,}32$~cm). Questo spiega perché l'errore
        grezzo cresce con la distanza (da $-0{,}5$~cm a 1~m fino a $+18{,}2$~cm a
        5~m) ed è, soprattutto, \textbf{correggibile con un solo coefficiente}:
        dividendo la lettura per 1,0381 l'errore residuo scende sotto i 5~cm a tutte
        le distanze.
\end{itemize}

\paragraph{Limite dichiarato sulla causa.} I dati raccolti \emph{non} permettono di
attribuire l'origine del $+3{,}81\,\%$. Può trattarsi della calibrazione interna del
modulo oppure del riferimento con cui sono state segnate le distanze sul pavimento.
Per discriminare servirebbe un riferimento di distanza indipendente, per esempio un
distanziometro laser. La distinzione è rilevante: nel primo caso il coefficiente
correttivo è una proprietà del modulo e vale in generale, nel secondo è un artefatto
del nostro setup.

\subsection{Dispersione entro il trial}

La dispersione della distanza \emph{entro} un singolo trial è risultata di
10--26~cm nello scenario di cammino sul posto, contro
$\mathbf{1{,}48 \pm 0{,}77}$~\textbf{cm} nello scenario da seduto immobile
(\cref{sec:persona-ferma}). Il confronto è informativo: la dispersione maggiore non
è rumore del sensore, è l'oscillazione reale del busto di chi cammina sul posto. Il
radar la sta misurando correttamente.

Nello scenario da seduto emerge una scomposizione utile: la variabilità \emph{fra}
trial della distanza media è di $\pm 2{,}38$~cm (media $229{,}78$~cm su una distanza
reale di 230~cm), mentre la stabilità \emph{entro} il trial è di $1{,}48$~cm. La
variabilità di come il soggetto si siede ($\sim$2,4~cm) è quindi maggiore di quella
del sensore ($\sim$1,5~cm): il fattore limitante dell'accuratezza in questo scenario
è il soggetto, non lo strumento.

\subsection{Decadimento dell'energia con la distanza}

L'energia media del bersaglio decade da 99 a 1~m fino a 28 a 5~m
(\cref{tab:distanza}). Il decadimento è \emph{molto} più lento di quanto
prevederebbe l'equazione del radar, che per un bersaglio puntiforme dà una
dipendenza $1/D^4$ della potenza ricevuta.

Va quindi evitata una tentazione interpretativa: \textbf{questo valore non è una
potenza e non va letto con l'equazione del radar}. È un indicatore normalizzato
0--100 prodotto da un'elaborazione interna non documentata, con ogni probabilità
compressiva. Se ne può usare la monotonia (più lontano $\to$ meno energia) ma non il
valore assoluto in termini fisici.

L'osservazione ha però un uso pratico: l'extrapolazione della tendenza porta a un
valore di circa 27 a 6~m, contro una soglia di rilevamento di 15 sul gate
corrispondente (\cref{tab:soglie}). Questo indica che il sensore avrebbe funzionato
anche a 6~m.

\paragraph{Perimetro sperimentale.} Il range coperto è \textbf{1--5~m} e non 6~m: è
un limite della stanza disponibile, non del sensore, e va dichiarato come tale. La
retta di regressione è costruita su cinque punti e 25 trial. Per lo scenario UPRISE
la limitazione è comunque irrilevante, perché la distanza di interesse è
\emph{inferiore al metro} --- la persona sotto il banco --- che è coperta dallo
scenario di \cref{sec:sotto-banco}.

\section{Latenze di rilevamento e di rilascio}
\label{sec:latenza}

\emph{Acquisizione in corso.}

\paragraph{Metodo previsto.} Poiché la ground truth è statica per file, la latenza
si misura con la convenzione dell'\textbf{evento a istante noto}: l'acquisizione
parte con la stanza vuota e il soggetto entra esattamente a $t = 10$~s, con
riferimento condiviso fra operatore e acquisizione. Lo script di analisi riceve
l'istante dell'evento come parametro (\texttt{-{}-event-time} per il rilevamento,
\texttt{-{}-release-time} per il rilascio) e calcola la latenza come distanza
temporale fra l'evento e la prima transizione stabile.

\paragraph{Attese e vincoli noti.} Tre elementi sono già stabiliti e vanno riportati
con i risultati.

\begin{itemize}[leftmargin=1.4em,itemsep=2pt]
  \item La latenza di rilascio del radar \textbf{non è un difetto}: è il
        \emph{no-one duration}, impostato a 5~s di fabbrica, che nella sessione di
        validazione ha prodotto una coda osservata di circa 10~s con un bersaglio
        fantasma in decadimento a 5,2~m (\cref{sec:coda}). Va misurata come tale e
        non conteggiata fra i falsi positivi.
  \item Il PIR ha un \textbf{block time di 2,5~s} dopo ogni rilascio, durante il
        quale è cieco per costruzione. Una latenza di rilevamento misurata subito
        dopo un rilascio è quindi contaminata da questo effetto e va segnalata.
  \item La direzione di ingresso conta, e in senso opposto per i due sensori: il PIR
        favorisce l'ingresso \emph{trasversale}, il radar quello \emph{radiale}
        (\cref{sec:fresnel}). La direzione va annotata per ogni trial e i due casi
        vanno riportati separatamente, altrimenti il confronto è ambiguo.
\end{itemize}

\section{Lo scenario UPRISE: persona sotto il banco}
\label{sec:sotto-banco}

\emph{Acquisizione in corso.}

È lo scenario che replica la condizione d'interesse del progetto: soggetto
accovacciato sotto un piano, a distanza inferiore al metro dal sensore. Rispetto
agli scenari precedenti presenta tre elementi nuovi, tutti già identificati:

\begin{enumerate}[leftmargin=1.8em,itemsep=3pt]
  \item \textbf{distanza ravvicinata}, quindi rischio di saturazione dei canali di
        energia (\cref{sec:saturazione});
  \item \textbf{canali stazionari per-gate non disponibili} sotto 1,5~m
        (\cref{sec:senergy-zero}): l'analisi dovrà usare i canali in movimento,
        e questa verifica è esplicitamente parte dello scenario;
  \item \textbf{presenza del piano dell'arredo} fra sensore e soggetto, a seconda
        del punto di montaggio. Nel banco UPRISE il piano è rinforzato con
        \emph{lamiera forata}: il metallo è il materiale peggiore per un radar
        (\cref{sec:ostacoli}), e la posizione di montaggio del sensore rispetto alla
        lamiera è una domanda aperta di progetto, da porre al relatore.
\end{enumerate}

\section{Penetrazione degli ostacoli}
\label{sec:ostacoli}

\emph{Acquisizione in corso.}

\paragraph{Metodo previsto.} Ripetizione dello scenario ``persona ferma'' con
l'interposizione di un pannello fra sensore e soggetto, per i materiali disponibili:
cartongesso, legno, vetro, plastica. Metrica: tasso di falsi negativi e riduzione
dell'energia del bersaglio rispetto al caso senza ostacolo, a parità di distanza.

\paragraph{Attese da letteratura.} Il radar a 24~GHz attraversa cartongesso, legno,
vetro e plastica con attenuazione contenuta, mentre è bloccato da metallo e da
cemento armato di spessore rilevante. Il PIR è invece bloccato da \emph{tutti} i
materiali elencati, vetro compreso: la finestra ottica 5--14~\textmu m del corpo
umano non attraversa il vetro comune. Quest'ultimo punto è un risultato atteso ma
importante per il confronto, perché rende il PIR inutilizzabile per qualunque
montaggio incassato o protetto.

\section{Il rilevamento del respiro}
\label{sec:respiro}

\subsection{Metodo}

Il metodo segue quanto descritto in \cref{sec:respiro-teoria}: si campiona a 5~Hz
un canale di energia per-gate in engineering mode, si applica la FFT alla serie
temporale e si cerca il picco in banda 0,1--0,5~Hz. La frequenza del picco,
moltiplicata per 60, dà la stima degli atti al minuto.

Il contributo metodologico di questa tesi è nella \textbf{selezione del canale}. Il
LD2410B mette a disposizione venti serie temporali di energia (due di bersaglio più
diciotto per-gate) e non tutte sono utilizzabili: alcune sono sature
(\cref{sec:saturazione}), altre identicamente nulle (\cref{sec:senergy-zero}), altre
piatte perché lontane dal bersaglio. La modalità \texttt{-{}-scan} dello script di
analisi esamina tutti i canali, scarta quelli saturi e quelli piatti secondo criteri
dichiarati, ordina i rimanenti per rapporto segnale-rumore e riporta la
\textbf{mediana delle stime concordi}. La concordanza fra canali indipendenti è
l'unico controllo di plausibilità disponibile in assenza di ground truth.

\subsection{Prima prova: soggetto a 40~cm}

\paragraph{Risultato.} Soggetto fermo a circa 40~cm, 91~s di acquisizione,
respirazione spontanea. \textbf{Sette canali di energia in movimento, indipendenti
fra loro, concordano su 0,264~Hz, cioè 15,8 atti al minuto}, con rapporto
segnale-rumore fino a 12,8 sul canale \texttt{menergy\_gate2}. Il valore è
fisiologicamente plausibile per un adulto a riposo.

\paragraph{Perché non è una misura valida.} Va detto con chiarezza, perché è una
distinzione che qualifica il lavoro: questa prova dimostra la \emph{fattibilità} ma
non è una misura utilizzabile come risultato. Le ragioni sono tre:
\emph{(i)} manca la ground truth, cioè non si conosce il ritmo respiratorio reale e
quindi non si può calcolare un errore; \emph{(ii)} l'89\,\% dei campioni sui canali
in movimento è saturo; \emph{(iii)} 40~cm non è una distanza rappresentativa di
alcuno scenario d'uso.

\paragraph{I canali stazionari sono inutilizzabili a questa distanza.} I gate 2 e 3
stazionari sono saturi al 100\,\%, i gate 0 e 1 sono identicamente nulli
(\cref{sec:senergy-zero}) e i gate 4--8 forniscono stime fra 6,6 e 9,2 atti al
minuto, in disaccordo fra loro e con il canale in movimento. Il valore 0,132~Hz che
appare su alcuni di questi canali è esattamente la metà di 0,264~Hz: si tratta con
ogni probabilità di una subarmonica, cioè di un artefatto dell'analisi, non di un
segnale.

\subsection{Seconda prova: soggetto a 2,31~m}

La stessa analisi è stata applicata ai cinque trial dello scenario ``persona ferma''
(\cref{sec:persona-ferma}), a 2,31~m e per 302~s ciascuno. Il quadro si ripete, ma
in condizioni migliori.

\begin{itemize}[leftmargin=1.4em,itemsep=3pt]
  \item I canali \textbf{in movimento} dei gate 2--8 concordano su circa 0,30~Hz,
        cioè \textbf{18,2 atti al minuto}, valore plausibile per un adulto seduto.
        Le stime per i cinque trial sono 18,2 / 18,9 / 18,8 / 20,0 / 20,2 atti al
        minuto, con 7--8 canali concordi per trial.
  \item Questi canali \textbf{non sono saturi} a questa distanza: la percentuale di
        campioni a fondo scala è fra 0 e 1\,\%. È una buona notizia di natura
        pratica, perché significa che la prova definitiva del respiro può essere
        svolta a distanza realistica senza accorgimenti geometrici particolari.
  \item I canali \textbf{stazionari} danno 6,5--7,8 atti al minuto, troppo pochi per
        un adulto a riposo, e i gate prossimi al bersaglio sono saturi (gate~3 al
        100\,\%, gate~4 al 97\,\%).
\end{itemize}

\paragraph{Conclusione operativa.} Due sessioni a distanze molto diverse (40~cm e
2,31~m) dicono la stessa cosa: \textbf{per il respiro vanno usati i canali in
movimento}. Il valore intorno a 7 atti al minuto dei canali stazionari è
verosimilmente un artefatto del filtraggio interno del canale e non respirazione ---
ma, in assenza di ground truth, questa resta un'ipotesi e va presentata come tale.
A deciderlo è la prova con metronomo descritta di seguito.

\subsection{La prova definitiva: respiro a ritmo imposto}
\label{sec:respiro-metronomo}

\emph{Acquisizione in corso.}

La misura valida per la tesi richiede una ground truth, e il modo più semplice di
ottenerla è imporre il ritmo: il soggetto respira seguendo un metronomo a
frequenza nota (per esempio 10, 15 e 20 atti al minuto), a distanza $\geq 1{,}5$~m,
per cinque trial per ritmo. Il confronto fra stima e ritmo imposto è la vera prova,
e permette di riportare un errore assoluto in atti al minuto invece di una
concordanza fra canali.

È inoltre necessario un \textbf{controllo negativo}: la stessa analisi applicata a
una registrazione di stanza vuota deve \emph{non} produrre un picco in banda
respiratoria. Senza questo controllo, un picco nello spettro non dimostra nulla, e
la banda 0,1--0,5~Hz è abbastanza larga da ospitare artefatti.

\section{Confronto fra i due moduli radar}
\label{sec:ld2420-risultati}

\emph{Attività parziale.}

Il LD2420 è stato collegato e verificato (\cref{sec:ld2420}): risponde ai movimenti,
riporta presenza e un valore di ``range'' in modalità ASCII a 115200 baud. In questa
modalità però non fornisce energia per-gate né distingue movimento e stazionario, e
l'unità del campo ``range'' non è documentata. Il confronto quantitativo con il
LD2410B richiede quindi la riconfigurazione del modulo in modalità binaria.

Alla data di questa stesura il confronto è limitato a un raffronto di specifiche
(\cref{tab:confronto-moduli}) e alla constatazione qualitativa che il LD2420
rileva presenza a distanze superiori. Il posizionamento dei due moduli resta quello
indicato dalle specifiche: il LD2410B è il modulo adatto al caso UPRISE (dettaglio
per-gate, respiro, distanze brevi), il LD2420 a una copertura di ambiente più ampia.

\section{Discussione}
\label{sec:discussione}

\subsection{Dove il mmWave vince}

Il quadro che emerge dai dati acquisiti è netto e si può riassumere in quattro
punti.

\begin{enumerate}[leftmargin=1.8em,itemsep=3pt]
  \item \textbf{Sulla persona ferma il divario non è graduale, è categorico}:
        $0{,}00\,\%$ contro $99{,}78\,\%$ di falsi negativi. Non si tratta di un
        sensore migliore di un altro, ma di un sensore che misura una grandezza che
        l'altro non ha accesso a misurare.
  \item \textbf{Il divario si estende alla persona in movimento sul posto}, che è il
        caso applicativo reale: $100\,\%$ di falsi negativi PIR oltre i 2~m contro
        $0\,\%$ del radar.
  \item \textbf{Il radar aggiunge informazione, non solo affidabilità}: la distanza
        è misurata con linearità $R^2 = 0{,}99965$ e un errore correggibile con un
        coefficiente; l'energia per-gate consente l'analisi del respiro e l'indice
        di vitalità (\cref{cap:vitalita}), cose che con un bit non sono
        rappresentabili.
  \item \textbf{Sui falsi positivi i due sensori sono equivalenti}: entrambi
        $\leq 0{,}43$ eventi/h su 7~h di osservazione. Il vantaggio del radar non si
        paga in affidabilità.
\end{enumerate}

\subsection{Dove il mmWave non vince}

Per equilibrio vanno riportati anche i punti a sfavore, che sono reali.

\begin{itemize}[leftmargin=1.4em,itemsep=3pt]
  \item \textbf{Consumo}: tre ordini di grandezza in più (80~mA contro 50~\textmu A).
        È il vincolo che, da solo, determina l'architettura raccomandata nel
        \cref{cap:conclusioni}, e il \cref{cap:consumi} lo quantifica.
  \item \textbf{Coda di presenza}: il radar mantiene la presenza per circa 10~s dopo
        l'uscita. Per un'applicazione di illuminazione automatica sarebbe un
        vantaggio, per una misura di occupazione istantanea è un ritardo.
  \item \textbf{Saturazione a distanza ravvicinata}: proprio nello scenario di
        interesse (sotto il banco) i canali più informativi rischiano di saturare,
        e i canali stazionari per-gate non sono disponibili sotto 1,5~m. Sono
        vincoli di progetto, non insormontabili, ma vanno gestiti.
  \item \textbf{Bersaglio singolo}: il LD2410B non dispone di un array di antenne e
        non distingue due persone alla stessa distanza. Non fornisce informazione
        angolare.
  \item \textbf{Complessità}: richiede UART a 256000 baud, gestione di un protocollo
        a frame, e --- per usarne i dati veri --- engineering mode e analisi
        numerica. Il PIR richiede un piedino digitale.
\end{itemize}

\subsection{Limiti del setup sperimentale}
\label{sec:limiti}

I limiti da dichiarare esplicitamente sono i seguenti.

\begin{description}[leftmargin=1.4em,style=nextline,itemsep=3pt]
  \item[Numero di soggetti] Le acquisizioni sono state svolte con un solo soggetto
    (indicato come S1 nel registro). La corporatura influisce sull'energia riflessa
    e, nel caso del PIR, sull'emissione infrarossa: la ripetizione degli scenari
    chiave con un secondo soggetto è la prima estensione da fare.
  \item[Ambiente] Stanza domestica, non un'aula scolastica: dimensioni, materiali,
    riflessioni multiple e presenza di arredi metallici sono diversi. Il
    trasferimento dei risultati al dimostratore reale va verificato.
  \item[Temperatura] Fra 25 e 29\,\celsius. Condizione sfavorevole al PIR in
    rilevamento: un confronto a 18--20\,\celsius\ è un'estensione a costo quasi
    nullo che rafforzerebbe il risultato, mostrando quanto della cecità sul fermo
    sia strutturale e quanto termico. La previsione, dalla fisica del principio, è
    che la cecità sul fermo resti totale a qualunque temperatura.
  \item[Range] 1--5~m per la distanza, per limite della stanza. Il caso UPRISE è
    sotto il metro e non è interessato.
  \item[Un solo esemplare per modello] Nessuna stima della variabilità fra esemplari.
    Il coefficiente di scala del $+3{,}81\,\%$ potrebbe essere specifico
    dell'esemplare in prova.
  \item[Assenza di strumentazione di misura] I consumi sono da datasheet
    (\cref{cap:consumi}) e le distanze sono riferite a marcature su pavimento, non a
    un distanziometro certificato.
\end{description}

Nessuno di questi limiti mette in discussione il risultato centrale, che è
robusto per la natura stessa del divario misurato: la cecità del PIR sulla persona
ferma è una conseguenza del principio fisico, e i dati la confermano nella
condizione di interesse con cinque ripetizioni indipendenti.

\chapter{Consumo energetico}
\label{cap:consumi}

\begin{quote}\small
Questo capitolo risponde all'obiettivo O4. L'analisi è dichiaratamente
\emph{a titolo informativo} e basata su dati di datasheet: non si dispone di
strumentazione di misura (né multimetro né resistenze di shunt), come concordato
con il relatore. Tutti i valori riportati sono quelli dichiarati dai produttori,
con la fonte indicata.
\end{quote}

\section{Consumo dei singoli componenti}

\subsection{Sensori}

\begin{table}[htbp]
\centering
\caption{Consumo dei sensori in esame, da datasheet.}
\label{tab:consumi-sensori}
\small
\begin{tabularx}{\textwidth}{@{}lllrX@{}}
\toprule
\textbf{Componente} & \textbf{Tensione} & \textbf{Corrente media} &
\textbf{Potenza} & \textbf{Fonte} \\
\midrule
\ldb & 5~V ($>$200~mA di capacità) & 80~mA & $\sim$400~mW &
  manuale Hi-Link V1.03~\cite{hilinkManualV103} \\
\ldc & 3,3~V (3,0--3,6~V) & $\sim$50~mA & $\sim$165~mW &
  datasheet Hi-Link / OpenELAB~\cite{ld2420openelab} \\
PIR \pir & 4,5--20~V & $<$0,05~mA & $\sim$0,3~mW &
  datasheet \pir~\cite{hcsr501} \\
\bottomrule
\end{tabularx}
\end{table}

Il rapporto fra i due estremi è immediato e va sottolineato perché è il dato
strutturale di tutto il capitolo: il radar consuma circa \textbf{1600 volte} la
corrente del PIR.

\paragraph{Nota su una fonte discordante.} Una fonte secondaria di uso
comune~\cite{components101} riporta per l'\pir\ una corrente di riposo di 65~mA. Si tratta con evidenza di un
refuso fra milliampere e microampere, incompatibile con l'architettura del modulo
(un BISS0001 alimentato da un LDO HT7133). Si è fatto fede al datasheet, che indica
$<$50~\textmu A. Il caso è riportato perché illustra la necessità della regola
adottata in tutta la tesi: risalire sempre al datasheet del produttore.

\subsection{Microcontrollore}

\begin{table}[htbp]
\centering
\caption{Consumo dell'ESP32-WROOM-32 per modalità operativa
(datasheet Espressif~\cite{esp32datasheet}).}
\label{tab:consumi-esp32}
\small
\begin{tabularx}{\textwidth}{@{}llX@{}}
\toprule
\textbf{Modalità} & \textbf{Corrente} & \textbf{Note} \\
\midrule
Attivo, WiFi in trasmissione & 160--260~mA (picco) & invio dati alla web UI \\
Attivo, WiFi in ricezione/idle & $\sim$95--100~mA & connesso, in ascolto \\
Modem-sleep (CPU attiva, WiFi spento) & 20--68~mA & dipende dalla frequenza di
  clock (80--240~MHz) \\
Light-sleep & $\sim$0,8~mA & \\
Deep-sleep & $\sim$0,01~mA (10~\textmu A) & solo dominio RTC attivo \\
\bottomrule
\end{tabularx}
\end{table}

\section{Consumo dei nodi completi}

Sommando microcontrollore e sensore si ottengono le configurazioni effettivamente
usate in questa tesi e quelle rilevanti per il progetto (\cref{tab:consumi-nodi}).

\begin{table}[htbp]
\centering
\caption{Consumo stimato del nodo completo nelle configurazioni di interesse.}
\label{tab:consumi-nodi}
\small
\begin{tabular}{@{}lrr@{}}
\toprule
\textbf{Configurazione} & \textbf{Corrente totale} & \textbf{Potenza (5~V)} \\
\midrule
ESP32 attivo + LD2410B + WiFi & $\sim$180--260~mA & $\sim$0,9--1,3~W \\
ESP32 attivo + LD2410B (solo logging seriale) & $\sim$110--150~mA & $\sim$0,55--0,75~W \\
ESP32 attivo + LD2420 & $\sim$80--120~mA & $\sim$0,4--0,6~W \\
ESP32 attivo + PIR & $\sim$30--70~mA & $\sim$0,15--0,35~W \\
ESP32 in deep-sleep + PIR (sveglia su interrupt) & $\mathbf{\sim 0{,}06}$~\textbf{mA} & $\sim$0,3~mW \\
\bottomrule
\end{tabular}
\end{table}

L'ultima riga è il punto chiave del capitolo. L'uscita digitale del PIR può
risvegliare l'ESP32 dal deep-sleep tramite interrupt su GPIO: il nodo può quindi
restare in attesa a costo praticamente nullo, con il PIR nel ruolo di
\emph{guardiano}. Il radar mmWave, invece, costa tre ordini di grandezza in più ed
è progettato per un'alimentazione fissa.

\section{Stime di autonomia a batteria}

Assumendo una cella litio-ione 18650 da 3000~mAh e un rendimento del regolatore di
circa l'85\,\%, si ottengono le autonomie di \cref{tab:autonomia}.

\begin{table}[htbp]
\centering
\caption{Autonomia stimata con una cella Li-ion 18650 da 3000~mAh e rendimento del
regolatore dell'85\,\%.}
\label{tab:autonomia}
\small
\begin{tabular}{@{}lrr@{}}
\toprule
\textbf{Scenario} & \textbf{Corrente media} & \textbf{Autonomia} \\
\midrule
Nodo mmWave sempre attivo (LD2410B + ESP32 + WiFi) & $\sim$200~mA & $\sim$13~h \\
Nodo mmWave attivo senza WiFi continuo & $\sim$130~mA & $\sim$20~h \\
Nodo PIR sempre attivo (ESP32 in modem-sleep) & $\sim$25~mA & $\sim$4 giorni \\
Nodo dormiente (deep-sleep + PIR come watchdog) & $\sim$0,06~mA & $\sim$4--5 anni$^{*}$ \\
\bottomrule
\multicolumn{3}{@{}l@{}}{\footnotesize $^{*}$ in questo regime il fattore
limitante non è il consumo ma l'autoscarica della cella.}
\end{tabular}
\end{table}

\section{Implicazione per UPRISE: l'architettura ibrida}
\label{sec:architettura-ibrida}

I numeri precedenti risolvono un'apparente contraddizione del progetto. Il sistema
integrato negli arredi deve restare in attesa \emph{per anni} in ``tempo di pace'' e
deve funzionare a batteria durante il blackout che segue il sisma: due requisiti
incompatibili con un sensore da 80~mA sempre acceso, ma anche con un sensore che non
vede le persone ferme.

La soluzione non è scegliere fra i due sensori, ma metterli in sequenza temporale.

\begin{enumerate}[leftmargin=1.8em,itemsep=4pt]
  \item \textbf{Tempo di pace} --- l'intero nodo in deep-sleep, con il PIR spento o
        usato come watchdog: consumo dell'ordine dei microampere, autonomia
        pluriennale limitata dall'autoscarica della batteria.
  \item \textbf{Trigger di ``modalità terremoto''} --- l'accelerometro sul gateway
        o il rilevamento del blackout risveglia l'ESP32. Il PIR può contribuire come
        sorgente di interrupt a costo nullo.
  \item \textbf{Emergenza} --- si accende il radar mmWave per il rilevamento fine:
        persona ferma sotto il banco, respiro, indice di vitalità. L'autonomia di
        13--20~ore copre la finestra critica iniziale delle operazioni di ricerca e
        soccorso.
  \item \textbf{Estensione della finestra} --- per coprire le prime 72~ore, che è
        l'orizzonte convenzionale del soccorso, sono disponibili due strade
        combinabili: una batteria maggiorata, oppure il \emph{duty-cycling} del
        radar. Un ciclo di 1~minuto acceso ogni 5 riduce il consumo medio di circa
        cinque volte, portando l'autonomia oltre le 72~ore. Il costo è una latenza
        di rilevamento fino a quattro minuti, che nello scenario di soccorso è
        accettabile: la domanda a cui il sistema risponde è ``c'è qualcuno sotto
        questo banco?'', non ``in quale istante si è mosso?''.
\end{enumerate}

\paragraph{Conclusione del capitolo.} L'analisi dei consumi porta a una
raccomandazione precisa, che è anche la risposta alla domanda implicita del titolo
della tesi: il PIR \textbf{non va sostituito}. Non è un concorrente del radar ma il
guardiano a basso costo che decide quando accenderlo. Il radar mmWave è lo strumento
di misura della fase di emergenza. Ciascuno dei due fa ciò che l'altro non può fare,
e i dati del \cref{cap:confronto} e di questo capitolo lo mostrano da due direzioni
opposte.

\todo{se il tempo lo consente: verificare sperimentalmente il consumo con un
misuratore USB in linea, anche a bassa precisione. Anche una misura al 10\,\%
di accuratezza confermerebbe l'ordine di grandezza delle tabelle e renderebbe il
capitolo più solido}

\chapter{Interfaccia web: visualizzazione e raccolta dati}
\label{cap:webui}

\begin{quote}\small
Questo capitolo risponde all'obiettivo O5: l'ESP32 pubblica i dati, un sito web li
visualizza in tempo reale, la stessa applicazione li salva con statistiche ed
esporta in CSV per l'analisi numerica. Alla data di questa stesura il progetto
architetturale e di dettaglio è completo; l'implementazione è in corso.
\end{quote}

\section{Requisiti e scelta architetturale}
\label{sec:webui-requisiti}

\subsection{Requisiti}

Dall'obiettivo O5 si ricavano cinque requisiti funzionali e due non funzionali.

\begin{description}[leftmargin=1.4em,style=nextline,itemsep=2pt]
  \item[R1 --- Pubblicazione] l'ESP32 rende disponibili i dati del sensore in rete.
  \item[R2 --- Tempo reale] la visualizzazione segue il sensore senza ritardi
    percepibili, alla stessa cadenza dell'acquisizione (5~Hz).
  \item[R3 --- Statistiche] la stessa applicazione calcola e mostra statistiche di
    sessione.
  \item[R4 --- Esportazione] i dati sono esportabili in CSV per l'analisi in Excel.
  \item[R5 --- Vitalità] l'interfaccia mostra anche l'indice di vitalità
    (\cref{cap:vitalita}).
  \item[N1 --- Funzionamento offline] vincolo derivato dal contesto UPRISE: in
    emergenza sismica non c'è connettività.
  \item[N2 --- Coerenza dei dati] il CSV esportato dal browser deve avere lo stesso
    formato di quello prodotto dalla catena seriale, in modo che esista
    \emph{un solo} formato dati in tutta la tesi.
\end{description}

\subsection{Opzioni considerate}

\begin{table}[htbp]
\centering
\caption{Opzioni architetturali valutate per la web UI.}
\label{tab:webui-opzioni}
\small
\begin{tabularx}{\textwidth}{@{}lXX@{}}
\toprule
\textbf{Architettura} & \textbf{Vantaggi} & \textbf{Svantaggi} \\
\midrule
\textbf{A. Sito servito dall'ESP32} (web server e WebSocket a bordo) &
  nessuna infrastruttura, funziona offline, dimostrazione autonoma, coerente con lo
  scenario di emergenza &
  RAM e flash limitate, massimo 2--3 client simultanei \\
B. ESP32 $\to$ MQTT $\to$ server con database &
  scalabile, storico illimitato, più vicino a una piattaforma reale &
  richiede broker e server sempre accesi: infrastruttura eccessiva per una tesi di
  sensing \\
C. Piattaforma cloud (ThingSpeak, Blynk) &
  pronta all'uso &
  non funziona offline, cadenza di campionamento limitata ($\sim$15~s), dati fuori
  controllo \\
\bottomrule
\end{tabularx}
\end{table}

\subsection{Scelta: tutto sull'ESP32}

È stata scelta l'opzione~A, per tre motivi.

\begin{enumerate}[leftmargin=1.8em,itemsep=3pt]
  \item \textbf{Coerenza con il contesto}: un nodo autonomo che serve la propria
        dashboard è la miniatura concettuale della piattaforma di monitoraggio del
        progetto, con le sue due modalità di funzionamento. L'opzione~C viola
        direttamente il requisito N1.
  \item \textbf{Il volume dei dati è piccolo}: circa venti valori a 5~Hz,
        dell'ordine di 2~kB/s, ben dentro le capacità dell'ESP32.
  \item \textbf{La persistenza a lungo termine non serve a bordo}: lo storico si
        accumula \emph{nel browser}, che ha memoria abbondante, ed è il browser a
        produrre il CSV. Questo elimina il bisogno di un database senza perdere
        funzionalità.
\end{enumerate}

\paragraph{Modalità WiFi doppia.} Il firmware supporta due modalità selezionabili:
\emph{Station}, in cui l'ESP32 si collega a una rete esistente (comodo in fase di
sviluppo), e \emph{Access Point}, in cui crea la propria rete e serve la pagina a
chi si collega. Se la connessione in modalità Station non riesce entro 15~s, parte
automaticamente l'Access Point. È la modalità che rende la dimostrazione possibile
in qualunque luogo, e quella che corrisponde allo scenario di emergenza.

\paragraph{Piano di riserva.} È stato progettato anche il percorso alternativo
(opzione~B: ESP32 $\to$ MQTT $\to$ backend FastAPI con SQLite $\to$ browser) nel
caso in cui il relatore intenda una piattaforma centralizzata anziché il nodo
autonomo. Il frontend è condiviso per circa l'85\,\% fra le due architetture, quindi
il cambio di rotta costerebbe due o tre giornate di lavoro. La scelta di riferimento
resta l'opzione~A.

\section{Progetto}
\label{sec:webui-progetto}

\subsection{Stack}

Lato firmware: \texttt{ESPAsyncWebServer} per servire la pagina, WebSocket per il
flusso dati, LittleFS per ospitare i file statici. Tutto risiede sull'ESP32 e nulla
è caricato da rete esterna, in ottemperanza al requisito N1.

Lato browser: pagina singola in HTML, CSS e JavaScript, senza framework e senza
dipendenze da CDN --- coerentemente con il vincolo di funzionamento offline.

\subsection{Protocollo dati}

Il messaggio WebSocket riprende le stesse grandezze del CSV, con l'aggiunta dei campi
calcolati a bordo:

\begin{lstlisting}[caption={Messaggio WebSocket inviato a 5~Hz.},
label={lst:ws},basicstyle=\ttfamily\scriptsize]
{
  "t": 123456,                        // millis() dell'ESP32
  "presence": 1,                      // radar_presence
  "moving": 1, "still": 0,            // stato del bersaglio
  "mdist": 145, "sdist": 0,           // distanze in cm
  "menergy": 72, "senergy": 0,        // energia del bersaglio 0-100
  "pir": 1,                           // sensore PIR (confronto dal vivo)
  "gates_m": [0,5,60,10,0,0,0,0,0],   // energia per-gate, movimento
  "gates_s": [0,0,45,8,0,0,0,0,0],    // energia per-gate, stazionario
  "vitality": 54,                     // indice di vitalita 0-100
  "vitality_class": "attivo"          // nessun_segno | vitalita_bassa |
                                      // moderato | attivo
}
\end{lstlisting}

La cadenza di invio è di \textbf{5~Hz}, identica al campionamento: nessun
sottocampionamento, così che il CSV esportato dal browser sia equivalente a quello
prodotto dalla catena seriale (requisito N2).

Il canale WebSocket gestisce la riconnessione automatica con backoff: se la
connessione cade, la pagina segnala lo stato e riprende senza perdere le statistiche
già accumulate.

\subsection{Struttura della pagina}

La pagina è organizzata in quattro aree, schematizzate in \cref{fig:webui-layout}.

\begin{figure}[htbp]
\centering
\begin{minipage}{0.95\textwidth}
\ttfamily\scriptsize
\begin{verbatim}
+-------------------------------------------------------------+
|  HEADER   * PRESENZA RILEVATA     [dist. 1.45 m]  [REC]     |
+--------------------------+----------------------------------+
|  A. STATO LIVE           |  B. INDICE DI VITALITA'          |
|  radar: MOVING @ 145 cm  |        +--- gauge 0-100 ---+     |
|  PIR:   ATTIVO           |        |        54         |     |
|  energia mov.:   72      |        |     "attivo"      |     |
|  energia still.: 30      |        +-------------------+     |
+--------------------------+----------------------------------+
|  C. GRAFICI                                                 |
|  C1: serie temporale energia moving/still + soglia           |
|  C2: barre per-gate (9 gate x moving/still) = "dov'e'"       |
|  C3: timeline presenza radar vs PIR (confronto visivo)       |
+-------------------------------------------------------------+
|  D. SESSIONE ED ESPORTAZIONE                                 |
|  scenario:[______] trial:[__] gt:[0|1]   <- metadati test    |
|  [Avvia] [Stop] [Scarica CSV] [Reset statistiche]            |
|  durata 12:34 | radar 97.2% | PIR 41.0% | eventi 3           |
+-------------------------------------------------------------+
\end{verbatim}
\end{minipage}
\caption{Struttura della pagina della dashboard. Il grafico C3 --- la timeline
affiancata di presenza radar e PIR --- è l'elemento che rende visivamente immediato
il risultato del \cref{cap:confronto}.}
\label{fig:webui-layout}
\end{figure}

Due scelte di progetto meritano una nota.

\begin{itemize}[leftmargin=1.4em,itemsep=3pt]
  \item Il grafico C3 mostra \textbf{radar e PIR affiancati sulla stessa base dei
        tempi}. È l'elemento più efficace dal punto di vista dimostrativo: un
        osservatore vede la riga del radar continua e quella del PIR a impulsi
        isolati mentre una persona sta ferma davanti al sensore. È il risultato
        della tesi reso visibile in tempo reale.
  \item L'area D contiene i \textbf{metadati sperimentali} (scenario, trial, ground
        truth) con gli stessi nomi usati dalla catena seriale. In questo modo il CSV
        scaricato dal browser è direttamente utilizzabile dagli stessi script di
        analisi, senza conversioni.
\end{itemize}

\subsection{Statistiche di sessione}

Le statistiche sono calcolate \emph{nel browser} e non sull'ESP32, per non caricare
il microcontrollore con un compito che non gli è proprio:

\begin{itemize}[leftmargin=1.4em,itemsep=1pt]
  \item durata della sessione e percentuale di campioni con presenza, separatamente
        per radar e PIR;
  \item numero di eventi di attivazione (fronti $0 \to 1$) per ciascun sensore, che
        è la stima dei falsi positivi in tempo reale;
  \item distanza minima, media e massima; energia media;
  \item tempo trascorso dall'ultima rilevazione --- l'informazione più rilevante
        nello scenario di soccorso;
  \item vitalità: valore corrente, media mobile, minimo e massimo di sessione.
\end{itemize}

\section{Piano di realizzazione e criteri di accettazione}
\label{sec:webui-piano}

Lo sviluppo è organizzato in cinque passi, ciascuno con un criterio di accettazione
verificabile. La suddivisione serve a garantire che ogni passo produca qualcosa di
utilizzabile anche se i successivi non vengono completati.

\begin{enumerate}[leftmargin=1.8em,itemsep=3pt]
  \item \textbf{Server e pagina statica} --- l'ESP32 serve una pagina da LittleFS in
        entrambe le modalità WiFi. \emph{Accettazione}: la pagina si apre da browser
        collegandosi all'Access Point del nodo.
  \item \textbf{Flusso WebSocket} --- invio del messaggio di
        \cref{lst:ws} a 5~Hz. \emph{Accettazione}: i valori sulla pagina cambiano
        insieme al sensore, senza ritardo percepibile e senza accumulo di latenza su
        dieci minuti di funzionamento.
  \item \textbf{Visualizzazione} --- stato live, i tre grafici, la gauge di vitalità.
        \emph{Accettazione}: le barre per-gate indicano il gate corretto quando il
        soggetto si sposta di 0,75~m.
  \item \textbf{Registrazione ed esportazione} --- accumulo nel browser ed
        esportazione CSV. \emph{Accettazione}: il CSV scaricato dal browser e quello
        prodotto dalla catena seriale, per la stessa sessione, sono
        \textbf{equivalenti}; gli script di analisi producono le stesse metriche su
        entrambi. Questo è il test che chiude il requisito N2.
  \item \textbf{Indice di vitalità a bordo} --- porting dell'algoritmo del
        \cref{cap:vitalita}. \emph{Accettazione}: l'indice calcolato sull'ESP32 e
        quello calcolato in Python sullo stesso CSV coincidono entro l'errore di
        arrotondamento.
\end{enumerate}

\section{Riferimenti di progetto}

Come riferimento di interfaccia è stato esaminato il \emph{LD2410
Configurator}~\cite{ld2410configurator}, un configuratore web open source che dialoga
con il modulo tramite Web Serial e Web Bluetooth. Da esso sono stati ripresi due
elementi: la rappresentazione a barre delle energie per-gate con la soglia
sovrapposta, che è il modo più leggibile di mostrare venti canali insieme, e
l'impostazione ``una pagina, nessuna installazione''. Non sono state riprese le
funzioni di configurazione dei parametri, che nel nostro caso sarebbero controproducenti:
le soglie devono restare fisse per tutta la campagna (\cref{sec:soglie}), e un
pulsante che le modifica dalla dashboard è un rischio, non una funzione.

\paragraph{Avvertenza sulla fonte.} Il LD2410 Configurator è un progetto di comunità
e non un prodotto Hi-Link: vale come riferimento di interfaccia e di implementazione,
non come fonte di dati tecnici sul modulo. Per quelli si è fatto riferimento
esclusivamente ai documenti ufficiali~\cite{hilinkManualV104,hilinkProtoV107}.

\todo{completare con gli screenshot della dashboard realizzata e con l'esito
del test di accettazione del passo 4 (equivalenza dei due CSV)}

\chapter{L'indice di vitalità}
\label{cap:vitalita}

\begin{quote}\small
Questo capitolo risponde all'obiettivo O6, formulato dal relatore come segue:
\emph{oltre alla presenza sì/no, creare un indice di vitalità --- un algoritmo che
valuti quanto una persona si muove: questa persona si muove 50 su 100, quindi è viva
o moderatamente viva}. È l'obiettivo più avanzato della tesi ed è anche la
dimostrazione costruttiva che i dati del radar abilitano ciò che con un bit non è
rappresentabile. Alla data di questa stesura l'algoritmo è specificato e il percorso
di taratura è definito; la taratura sui dati è in corso.
\end{quote}

\section{Motivazione: dalla presenza al triage}
\label{sec:vitalita-motivazione}

Nel contesto del progetto, la differenza fra ``presenza'' e ``vitalità'' è la
differenza fra due domande distinte dei soccorritori:

\begin{enumerate}[leftmargin=1.8em,itemsep=2pt]
  \item \emph{C'è qualcuno sotto il banco?} --- è la domanda di presenza, cui
        rispondono gli obiettivi O1--O3.
  \item \emph{In che condizioni si trova?} --- si muove attivamente, quindi è
        cosciente; compie solo micro-movimenti, quindi è ferito o intrappolato ma
        vivo; oppure non dà segni rilevabili.
\end{enumerate}

Il PIR non può nemmeno porsi la seconda domanda: la sua uscita è un bit, e come
mostrato in \cref{sec:interpretazione-pir} quel bit non misura nemmeno la durata del
movimento. Il radar mmWave sì, perché espone l'\textbf{energia riflessa per gate di
distanza}: una misura continua di \emph{quanto} si muove ciò che riflette, non solo
del fatto che si muova.

L'indice di vitalità è la sintesi di questa informazione in un numero fra 0 e 100,
leggibile da un operatore non tecnico, accompagnato da una classificazione a quattro
livelli.

\section{Dati di ingresso}
\label{sec:vitalita-ingressi}

L'algoritmo lavora sui campioni a 5~Hz prodotti dal logger
(\cref{sec:piattaforma}), usando i canali di \cref{tab:vitalita-ingressi}.

\begin{table}[htbp]
\centering
\caption{Segnali di ingresso dell'indice di vitalità.}
\label{tab:vitalita-ingressi}
\small
\begin{tabularx}{\textwidth}{@{}llX@{}}
\toprule
\textbf{Segnale} & \textbf{Colonna CSV} & \textbf{Che cosa misura} \\
\midrule
Energia del bersaglio in movimento & \texttt{moving\_energy} &
  intensità del movimento rilevato (0--100) \\
Energia per-gate, movimento & \texttt{menergy\_gate0..8} &
  dove avviene il movimento \\
Energia per-gate, stazionario & \texttt{senergy\_gate0..8} &
  riflessione del corpo fermo; oscilla con il respiro \\
Stato di presenza & \texttt{radar\_presence} &
  condizione di guardia: se nulla, l'indice è forzato a zero \\
\bottomrule
\end{tabularx}
\end{table}

L'osservazione su cui poggia l'algoritmo è la seguente: una persona \textbf{viva ma
immobile} ha energia in movimento prossima a zero, ma l'energia del suo gate
\emph{oscilla} con il respiro nella banda 0,1--0,5~Hz. Una stanza vuota ha energia
costante, a meno del rumore di fondo. La \textbf{variabilità} dell'energia è dunque
il segnale di vita minimo, quello che sta sotto il movimento.

\paragraph{Vincolo importante ereditato dal \cref{cap:moduli}.} La formulazione
iniziale dell'algoritmo usava l'energia \emph{stazionaria} del gate attivo come
sorgente del termine respiratorio. I risultati di \cref{sec:senergy-zero} e
\cref{sec:respiro} impongono una correzione: sotto 1,5~m i canali stazionari per-gate
sono strutturalmente nulli, e a tutte le distanze provate i canali stazionari
risultano saturi o discordanti mentre i canali \emph{in movimento} forniscono stime
concordi e non sature. L'implementazione deve quindi selezionare la sorgente del
termine respiratorio fra i canali in movimento, o comunque scartare i canali saturi
e quelli identicamente nulli con criterio esplicito. È un esempio di come la fase di
caratterizzazione abbia modificato il progetto dell'algoritmo, e non solo confermato
delle attese.

\section{Algoritmo}
\label{sec:vitalita-algoritmo}

L'algoritmo ha due componenti, aggiornate a ogni campione, e una fusione.

\begin{align}
g^{*} &= \arg\max_i \; e_i && \text{gate attivo} \label{eq:gate-attivo}\\
m_k &= \max\!\left(E_{\text{mov}},\; e^{\text{mov}}_{g^{*}}\right) && \text{movimento istantaneo} \\
M_k &= \alpha_m \, m_k + (1-\alpha_m)\, M_{k-1} && \text{EWMA rapida, } \alpha_m = 0{,}1 \\
d_k &= \left| e_{g^{*}}[k] - e_{g^{*}}[k-1] \right| && \text{variazione tick a tick} \\
R_k &= \alpha_r \, d_k + (1-\alpha_r)\, R_{k-1} && \text{EWMA lenta, } \alpha_r = 0{,}02 \\
V_k &= \mathrm{clamp}\!\left(M_k + k_r R_k,\; 0,\; 100\right) && \text{indice di vitalità}
\label{eq:vitalita}
\end{align}

con la condizione di guardia $V_k = 0$ se \texttt{radar\_presence} $= 0$.

\subsection{Razionale delle scelte}

Ogni scelta ha una motivazione, ed è utile esplicitarla perché l'algoritmo deve
girare identico in Python e su microcontrollore.

\begin{description}[leftmargin=1.4em,style=nextline,itemsep=4pt]
  \item[Media esponenziale e non finestra mobile] La EWMA costa $O(1)$ in memoria
    (due variabili in virgola mobile) e gira senza modifiche sull'ESP32; una finestra
    mobile richiederebbe un buffer circolare e darebbe risultati equivalenti. Poiché
    l'algoritmo deve essere lo stesso in Python e in C++ per essere verificabile
    (passo~5 di \cref{sec:webui-piano}), la formulazione più semplice è quella
    giusta.
  \item[Due costanti di tempo diverse] Il movimento deve reagire in circa 2~s, che a
    5~Hz corrisponde a $\alpha_m = 0{,}1$; il respiro è un segnale lento e va
    integrato su circa 10~s ($\alpha_r = 0{,}02$) per emergere dal rumore. Un'unica
    costante non può servire entrambi gli scopi.
  \item[Massimo fra energia del bersaglio e del gate] Le due grandezze talvolta
    divergono, per effetto dei filtri interni del modulo. Prendere il massimo rende
    l'indice \emph{conservativo verso il falso ``nessun segno''}, che è l'errore più
    grave nel dominio applicativo: dichiarare vuoto un banco sotto cui c'è qualcuno.
  \item[Il coefficiente $k_r$] La variazione $d_k$ è numericamente piccola, di
    qualche unità di energia; $k_r$ la riporta nella scala 10--30 attesa per la classe
    ``vitalità bassa''. È il parametro principale da tarare, con un intervallo
    iniziale plausibile fra 5 e 15.
\end{description}

\section{Classificazione}
\label{sec:vitalita-classi}

\begin{table}[htbp]
\centering
\caption{Classi dell'indice di vitalità e loro significato operativo. Le soglie
riportate sono iniziali e vanno tarate sui dati (\cref{sec:vitalita-taratura}).}
\label{tab:vitalita-classi}
\small
\begin{tabularx}{\textwidth}{@{}llX@{}}
\toprule
\textbf{Indice} & \textbf{Classe} & \textbf{Significato operativo (triage)} \\
\midrule
0--9 & \texttt{nessun\_segno} & nessuna presenza, o nessuna variazione rilevabile \\
10--39 & \texttt{vitalita\_bassa} & vivo: respiro e micro-movimenti, non si muove \\
40--69 & \texttt{moderato} & movimenti limitati ma attivi \\
70--100 & \texttt{attivo} & movimento pieno, persona reattiva \\
\bottomrule
\end{tabularx}
\end{table}

\paragraph{Avvertenza da riportare anche nella tesi finale.} Le classi sono un
supporto informativo al triage e \textbf{non costituiscono una diagnosi medica}. In
particolare \texttt{nessun\_segno} significa ``il sensore non rileva variazioni'', non
``persona deceduta'': la persona può essere fuori portata, schermata da materiale
metallico, oppure priva di sensi ma viva. La formulazione dell'etichetta, in
un'interfaccia destinata a soccorritori, è essa stessa una scelta di progetto con
implicazioni: un'etichetta che suggerisca una conclusione clinica sarebbe
inappropriata.

\section{Percorso di sviluppo e taratura}
\label{sec:vitalita-taratura}

La regola metodologica adottata è: \textbf{l'algoritmo si sviluppa su PC, sui file
CSV, dove si può iterare in secondi}. Il porting su ESP32 avviene solo a soglie
validate. Il percorso è il seguente.

\begin{enumerate}[leftmargin=1.8em,itemsep=4pt]
  \item \textbf{Raccolta dei quattro scenari}, cinque trial ciascuno:
        \texttt{vitalita\_0} (stanza vuota), \texttt{vitalita\_respiro} (soggetto
        immobile), \texttt{vitalita\_micro} (micro-movimenti: mano, testa, busto),
        \texttt{vitalita\_attivo} (movimento pieno). Distanza $\geq 1{,}5$~m per
        evitare la saturazione (\cref{sec:saturazione}).
  \item \textbf{Prototipo Python} (\texttt{analisi/vitalita\_proto.py}): implementa
        le equazioni \eqref{eq:gate-attivo}--\eqref{eq:vitalita} e produce la serie
        $V(t)$ per ogni trial.
  \item \textbf{Taratura} di $\alpha_m$, $\alpha_r$, $k_r$ e delle tre soglie in modo
        che i quattro scenari cadano nelle quattro classi attese. Il metodo scelto è
        semplice e difendibile: si osservano le distribuzioni di $V$ per scenario
        (diagrammi a scatola) e si collocano le soglie negli intervalli vuoti fra le
        distribuzioni. Se le distribuzioni si sovrappongono, la sovrapposizione è
        essa stessa un risultato da riportare, e indica che quei due scenari non sono
        separabili con questo indice.
  \item \textbf{Validazione su trial non usati per la taratura}: taratura su
        T01--T03, validazione su T04--T05, con matrice di confusione
        scenario/classe. La metrica finale è la percentuale di campioni classificati
        nella classe attesa, per scenario.
  \item \textbf{Porting} in \texttt{vitality.h} con le stesse costanti, e verifica
        dal vivo nella dashboard (passo~5 di \cref{sec:webui-piano}).
\end{enumerate}

Il passo~4 fornisce il numero da riportare nella tesi: \emph{l'indice classifica
correttamente l'$X\,\%$ dei campioni su scenari mai visti in fase di taratura}. La
separazione fra trial di taratura e trial di validazione è ciò che distingue questo
risultato da una descrizione dei dati su cui l'algoritmo è stato costruito.

\todo{completare con i risultati della taratura: boxplot per scenario, valori
scelti dei parametri, matrice di confusione sui trial di validazione}

\section{Casi limite e comportamenti attesi}
\label{sec:vitalita-limiti}

\begin{table}[htbp]
\centering
\caption{Casi limite dell'indice di vitalità e comportamento previsto.}
\label{tab:vitalita-casilimite}
\small
\begin{tabularx}{\textwidth}{@{}XXX@{}}
\toprule
\textbf{Situazione} & \textbf{Comportamento atteso} & \textbf{Note} \\
\midrule
La persona esce dal campo & indice a 0 entro il timeout del radar ($\sim$5~s) &
  forzatura sulla condizione di presenza \\
La persona si addormenta & da \texttt{attivo} a \texttt{vitalita\_bassa} in
  30--60~s & le medie esponenziali decadono con le loro costanti \\
Ventilatore o tenda in movimento & possibile falso \texttt{moderato} &
  limite da dichiarare; mitigazione tramite le soglie per-gate del radar \\
Due persone & indice riferito al gate dominante &
  coerente con il limite mono-bersaglio del LD2410B \\
Persona oltre 4--5~m & respiro non più rilevabile, quindi sottostima &
  va documentata la distanza massima utile per la classe
  \texttt{vitalita\_bassa} \\
Persona a meno di 1,5~m & canali stazionari per-gate non disponibili &
  usare i canali in movimento (\cref{sec:senergy-zero}) \\
\bottomrule
\end{tabularx}
\end{table}

\section{Estensioni possibili}

Tre estensioni sono state individuate ma non sono necessarie per l'obiettivo, e
vanno affrontate solo se il tempo lo consente.

\begin{enumerate}[leftmargin=1.8em,itemsep=2pt]
  \item Sostituire il termine respiratorio a media esponenziale con la
        \textbf{potenza in banda 0,1--0,5~Hz} calcolata su una finestra scorrevole di
        30~s, cioè una FFT ridotta eseguita a bordo. È più robusta agli artefatti ma
        computazionalmente più costosa, e richiede un buffer: sarebbe una versione~2.
  \item \textbf{Isteresi sulle classi}: una classe cambia solo dopo $N$ secondi di
        stabilità, per evitare lo sfarfallio dell'etichetta nell'interfaccia.
  \item \textbf{Indicatore di confidenza} affiancato all'indice, funzione della
        stabilità del gate attivo: un indice calcolato su un gate che cambia
        continuamente è meno affidabile e sarebbe onesto dichiararlo.
\end{enumerate}

\chapter{Conclusioni e sviluppi futuri}
\label{cap:conclusioni}

\section{Risposta alla domanda di partenza}

La domanda con cui si apre questa tesi era: \emph{il radar mmWave può sostituire o
affiancare il sensore PIR nel nodo DIPME-DEVICE del progetto UPRISE/SAFE?}

La risposta, sulla base dei dati raccolti, è \textbf{affiancare, non sostituire}, e
si articola in tre affermazioni sostenute da misure.

\paragraph{1. Sul caso d'uso che interessa il progetto, il PIR non funziona --- e
non per difetto ma per costruzione.} Con un soggetto immobile a 2,30~m, su cinque
trial da cinque minuti, il PIR ha mancato la presenza nel
$99{,}78 \pm 0{,}49\,\%$ dei campioni, con zero rilevamenti in quattro trial su
cinque; il radar l'ha rilevata nel $100\,\%$ dei 7554 campioni
(\cref{sec:persona-ferma}). Il risultato più significativo è però un altro: il PIR
fallisce anche con il soggetto \emph{in movimento sul posto}, con
$80{,}5\,\%$ di falsi negativi a 1~m e $100\,\%$ oltre i 2~m
(\cref{sec:pir-movimento}). Poiché una persona intrappolata si muove sul posto e non
attraversa la stanza, l'ipotesi implicita ``il PIR non vede chi è fermo ma vede chi
si muove'' è, nello scenario di interesse, falsa. Questa è la conclusione più
rilevante per il progetto.

\paragraph{2. Il radar non è solo più affidabile: fornisce grandezze diverse.} La
distanza è misurata con linearità eccellente ($R^2 = 0{,}99965$ su 25 trial fra 1 e
5~m) e con un errore riducibile sotto i 5~cm applicando un singolo coefficiente
correttivo (\cref{sec:distanza}). Le venti serie di energia dell'engineering mode
consentono di estrarre il ritmo respiratorio di un soggetto immobile, con
concordanza fra canali indipendenti (\cref{sec:respiro}), e sono la base
dell'indice di vitalità (\cref{cap:vitalita}). Nessuna di queste grandezze è
rappresentabile con l'uscita a un bit di un PIR. E il vantaggio non si paga in
affidabilità: su 7~ore di stanza vuota entrambi i sensori hanno prodotto zero falsi
positivi, con limite superiore al 95\,\% di $\leq 0{,}43$ eventi/h
(\cref{sec:falsi-positivi}).

\paragraph{3. Il consumo esclude la sostituzione secca e indica l'architettura.} Il
radar assorbe 80~mA contro i meno di 50~\textmu A del PIR: tre ordini di grandezza
(\cref{cap:consumi}). Un nodo integrato in un arredo deve restare in attesa per anni
a batteria, e con il radar sempre acceso l'autonomia si misura in ore. La soluzione
non è scegliere fra i due sensori ma metterli in sequenza temporale, secondo
l'architettura ibrida di \cref{sec:architettura-ibrida}: il PIR come guardiano a
microampere in ``tempo di pace'', il radar mmWave come strumento di misura durante la
finestra di emergenza, con \emph{duty-cycling} per estendere l'autonomia oltre le 72
ore. Il PIR non è un concorrente del radar: è ciò che decide quando accenderlo.

\medskip
\noindent
La raccomandazione operativa per il progetto è quindi di \textbf{aggiungere} un
modulo mmWave al DIPME-DEVICE, mantenendo il PIR nel ruolo di sorgente di risveglio,
e di attivare il radar sul trigger di modalità terremoto.

\section{Contributi consolidati}

Riassumendo quanto prodotto:

\begin{itemize}[leftmargin=1.4em,itemsep=3pt]
  \item una piattaforma di misura riproducibile su ESP32 che campiona radar e PIR in
        modo sincrono a 5,00~Hz esatti con accesso ai dati per-gate, con un formato
        CSV unico condiviso da firmware, acquisizione, analisi e interfaccia web;
  \item la caratterizzazione quantitativa, con cinque ripetizioni per scenario, del
        divario PIR/mmWave sulla persona ferma e sulla persona in movimento sul
        posto;
  \item la calibrazione della misura di distanza del LD2410B e l'individuazione di un
        errore di scala sistematico correggibile con un coefficiente;
  \item la dimostrazione di misurabilità del respiro con hardware consumer, con
        l'individuazione documentata di quali canali siano effettivamente
        utilizzabili e perché gli altri non lo siano;
  \item la specifica di un indice di vitalità implementabile su microcontrollore, con
        un percorso di taratura e validazione su trial separati;
  \item l'analisi dei consumi e l'argomentazione dell'architettura ibrida.
\end{itemize}

\subsection{Risultati metodologici collaterali}

Tre esiti non previsti meritano una menzione, perché sono il tipo di informazione
che raramente entra nella documentazione ma che ha valore per chi ripete il lavoro.

\begin{enumerate}[leftmargin=1.8em,itemsep=3pt]
  \item \textbf{Non fidarsi degli strumenti di configurazione da PC per i valori
        numerici.} Lo strumento ufficiale \texttt{LD2410 Tool} riporta sensibilità
        pari a 100 su tutti i gate; la lettura via UART restituisce la scala di
        fabbrica decrescente, coerente con il protocollo e con il comportamento
        osservato (\cref{sec:identita}). Per i parametri va usata la lettura diretta
        dal modulo.
  \item \textbf{Fare riferimento alla tabella di piedinatura del datasheet, non
        all'ordine apparente dei fili.} Un'inversione dei due fili della seriale ha
        prodotto tre settimane di lavoro nella convinzione che il modulo fosse
        guasto (\cref{sec:cablaggio}).
  \item \textbf{Verificare la qualità dei dati prima di analizzarli.} La saturazione
        dei canali di energia, i canali stazionari strutturalmente nulli e la coda di
        presenza (\cref{sec:insidie}) sono tutti fenomeni che, non riconosciuti,
        producono conclusioni sbagliate anziché errori evidenti: un canale saturo non
        segnala il proprio stato, restituisce semplicemente numeri privi di
        informazione. Da qui la scelta di eseguire un controllo di qualità
        automatico su ogni file prima di usarlo.
\end{enumerate}

\section{Lavoro rimanente}

Per completezza si elencano le attività pianificate e non ancora completate alla
data di questa stesura, con la sezione in cui sono collocate:

\begin{itemize}[leftmargin=1.4em,itemsep=2pt]
  \item latenze di rilevamento e di rilascio con la convenzione dell'evento a
        istante noto (\cref{sec:latenza});
  \item scenario ``persona sotto il banco'' a distanza inferiore al metro
        (\cref{sec:sotto-banco});
  \item penetrazione degli ostacoli: cartongesso, legno, vetro, plastica
        (\cref{sec:ostacoli});
  \item prova del respiro con ritmo imposto da metronomo e controllo negativo su
        stanza vuota (\cref{sec:respiro-metronomo});
  \item riconfigurazione del LD2420 in modalità binaria per il confronto quantitativo
        fra i due radar (\cref{sec:ld2420-risultati});
  \item implementazione della dashboard e test di equivalenza dei due CSV
        (\cref{sec:webui-piano});
  \item taratura e validazione dell'indice di vitalità
        (\cref{sec:vitalita-taratura}).
\end{itemize}

\section{Sviluppi futuri}

\paragraph{Verso il dimostratore reale.} La campagna è stata svolta in ambiente
domestico con un solo soggetto. Il passo successivo naturale è la ripetizione degli
scenari chiave nei dimostratori del progetto --- il Polo Lodovici di Informatica e
l'IIS Fermi Sacconi Ceci di Ascoli Piceno --- dove dimensioni dell'ambiente,
materiali, riflessioni multiple e presenza di più persone sono realistici. Le
domande aperte sono due: il tasso di falsi positivi in un'aula con arredi metallici,
e il comportamento con più persone contemporaneamente presenti, che è il limite
strutturale del LD2410B (\cref{sec:discussione}).

\paragraph{Il montaggio sulla lamiera forata.} È la questione di progetto più
concreta e va posta ai progettisti dell'arredo. Il piano del banco è rinforzato con
lamiera forata antisfondamento, e il metallo è il materiale peggiore per un radar a
24~GHz. Tre configurazioni sono da valutare sperimentalmente: sensore montato
\emph{sotto} il piano e quindi rivolto verso il volume di rifugio (ipotesi
preferibile); sensore montato \emph{sopra} il piano, che richiede di verificare se il
diametro e il passo dei fori consentono la trasmissione a 24~GHz --- una griglia
metallica con fori molto più piccoli della lunghezza d'onda di 12,5~mm si comporta
come uno schermo continuo, quindi la geometria della perforazione è determinante;
sensore integrato nel telaio laterale. Questa verifica è un'estensione di
\cref{sec:ostacoli} con il materiale reale.

\paragraph{Selettività fra banchi adiacenti.} Nel banco UPRISE i banchi sono
strutturalmente connessi fra loro, e un radar con portata di 6~m vede diversi banchi.
Il problema si affronta limitando il gate massimo via UART, funzione già
implementata e verificata nel corredo software di questa tesi. La misura da fare è
il tasso di attribuzione corretta con due volumi di rifugio adiacenti, che è
l'informazione necessaria a una mappa di calore utile ai soccorritori: sapere che
c'è qualcuno in un'aula è molto meno utile che sapere sotto quale banco.

\paragraph{Integrazione con il nodo DIPME e la rete LoRa.} L'indice di vitalità è
progettato per essere calcolato a bordo e occupa pochi byte: è esattamente il tipo di
informazione trasmissibile su LoRa, dove il vincolo è la dimensione del messaggio e
il duty cycle. Un messaggio che riporti presenza, distanza, classe di vitalità e
istante dell'ultima rilevazione sta in pochi byte ed è direttamente rappresentabile
sulla mappa di calore dei soccorritori. L'integrazione con il DIPME-DEVICE e con il
Digital Twin del progetto~\cite{callisto2023dt} è il proseguimento naturale del
lavoro.

\paragraph{Analisi in frequenza a bordo.} L'estensione dell'indice di vitalità con
una FFT calcolata sull'ESP32 su finestra scorrevole di 30~s renderebbe il termine
respiratorio più robusto agli artefatti rispetto alla media esponenziale attuale, a
costo di un buffer e di qualche millisecondo di calcolo per finestra: entrambi
disponibili sull'hardware in uso.

\paragraph{Confronto diretto con l'UWB del DIPME.} Il nodo esistente monta già un
sensore UWB. Un confronto sperimentale a tre --- PIR, mmWave, UWB --- sullo stesso
scenario e con la stessa metodologia sarebbe il completamento naturale di questo
lavoro, e trasformerebbe l'argomentazione di \cref{sec:uwb}, oggi basata su
datasheet e letteratura, in un risultato misurato.


\appendix
\chapter{Protocollo seriale del LD2410B}
\label{app:protocollo}

Riferimento: \emph{LD2410B Serial communication protocol}
V1.07~\cite{hilinkProtoV107}. Questa appendice riporta gli elementi effettivamente
usati dal firmware realizzato, per rendere il lavoro riproducibile senza dover
ricostruire il protocollo dal documento originale.

\section{Parametri della porta seriale}

\begin{center}\small
\begin{tabular}{@{}ll@{}}
\toprule
Baud rate & 256000 (fisso, non modificabile senza comando \texttt{0x00A1}) \\
Formato & 8 bit di dati, nessuna parità, 1 bit di stop (8N1) \\
Livelli & 3,3~V TTL \\
\bottomrule
\end{tabular}
\end{center}

Il baud di 256000 richiede una UART hardware: sull'ESP32 si usa \texttt{Serial2}
(GPIO25 come RX, GPIO26 come TX nella nostra configurazione). Un'implementazione a
software non regge questa velocità.

\section{Struttura dei frame}

\subsection{Frame di comando (dall'host al modulo)}

\begin{center}\small
\begin{tabular}{@{}lll@{}}
\toprule
\textbf{Campo} & \textbf{Lunghezza} & \textbf{Valore} \\
\midrule
Intestazione & 4 byte & \texttt{FD FC FB FA} \\
Lunghezza dati & 2 byte & little-endian \\
Parola di comando & 2 byte & little-endian \\
Valore del comando & variabile & dipende dal comando \\
Chiusura & 4 byte & \texttt{04 03 02 01} \\
\bottomrule
\end{tabular}
\end{center}

Il modulo risponde con un frame della stessa struttura, in cui la parola di comando
ha il bit più significativo posto a uno, seguita da uno stato (zero = successo) e
dagli eventuali dati di risposta.

\subsection{Frame di dati (dal modulo all'host)}

\begin{center}\small
\begin{tabular}{@{}lll@{}}
\toprule
\textbf{Campo} & \textbf{Lunghezza} & \textbf{Valore} \\
\midrule
Intestazione & 4 byte & \texttt{F4 F3 F2 F1} \\
Lunghezza dati & 2 byte & little-endian \\
Tipo di dato & 1 byte & \texttt{0x01} = engineering mode, \texttt{0x02} = normale \\
Testa & 1 byte & \texttt{0xAA} \\
Stato del bersaglio & 1 byte & maschera: \texttt{0x01} in movimento, \texttt{0x02} stazionario \\
Distanza in movimento & 2 byte & cm, little-endian \\
Energia in movimento & 1 byte & 0--100 \\
Distanza stazionaria & 2 byte & cm, little-endian \\
Energia stazionaria & 1 byte & 0--100 \\
Distanza di rilevamento & 2 byte & cm, little-endian \\
\multicolumn{3}{@{}l@{}}{\emph{solo in engineering mode:}} \\
Gate massimo movimento & 1 byte & \\
Gate massimo stazionario & 1 byte & \\
Energie per-gate movimento & 9 byte & gate 0--8, valori 0--100 \\
Energie per-gate stazionario & 9 byte & gate 0--8, valori 0--100 \\
Sensore di luce & 1 byte & 0--255 \\
Stato del pin OUT & 1 byte & \\
\multicolumn{3}{@{}l@{}}{} \\
Coda & 1 byte & \texttt{0x55} \\
Controllo & 1 byte & \texttt{0x00} \\
Chiusura & 4 byte & \texttt{F8 F7 F6 F5} \\
\bottomrule
\end{tabular}
\end{center}

\section{Comandi utilizzati}

\begin{center}\small
\begin{tabularx}{\textwidth}{@{}llX@{}}
\toprule
\textbf{Parola} & \textbf{Funzione} & \textbf{Note d'uso} \\
\midrule
\texttt{0x00FF} & entra in configurazione & obbligatorio prima di ogni altro
  comando; valore \texttt{0x0001} \\
\texttt{0x00FE} & esce dalla configurazione & il modulo riprende a inviare frame di
  dati \\
\texttt{0x0062} & attiva engineering mode & va inviato dentro la modalità
  configurazione \\
\texttt{0x0063} & disattiva engineering mode & \\
\texttt{0x00A0} & legge la versione firmware & restituisce tipo e versione
  (nel nostro caso 2.44.25070917) \\
\texttt{0x00A5} & legge il MAC Bluetooth & \\
\texttt{0x0061} & legge i parametri correnti & restituisce risoluzione, gate massimi,
  timeout e le 9+9 soglie; \textbf{è la lettura da usare per i parametri}, non lo
  strumento da PC \\
\texttt{0x0060} & imposta i parametri & usato per limitare il gate massimo nella prova
  di selettività spaziale \\
\texttt{0x00AA} & imposta la risoluzione di distanza & 0,75~m oppure 0,20~m per gate \\
\texttt{0x000B} & auto-calibrazione del rumore di fondo & disponibile dal firmware
  V2.44~\cite{hilinkV244} \\
\bottomrule
\end{tabularx}
\end{center}

\section{Regola di rilevamento}

Dal protocollo V1.07, \S1.2.2: il bersaglio in un gate è riconosciuto
\textbf{solo se l'energia del gate supera la soglia configurata per quel gate}. Da
questa regola derivano due fatti osservati sperimentalmente:

\begin{itemize}[leftmargin=1.4em,itemsep=2pt]
  \item i canali stazionari dei gate 0 e 1, che hanno soglia di fabbrica pari a~0 e
        energia riportata sempre nulla, sono disattivati per costruzione
        (\cref{sec:senergy-zero});
  \item le soglie influiscono sulla \emph{decisione} di rilevamento, non sul valore
        di energia riportato: non possono quindi risolvere la saturazione
        (\cref{sec:saturazione}).
\end{itemize}

\section{Auto-calibrazione del firmware V2.44}

Il documento delle novità del firmware V2.44~\cite{hilinkV244} descrive il
rilevamento automatico del rumore di fondo, che tara le soglie di movimento e
stazionarie. La procedura richiede 10~s per uscire dal campo del sensore più 60~s di
misura, per un totale di 70~s da trascorrere fuori dalla portata del modulo, e
l'applicazione mobile in versione adeguata (Android $\geq$ V1.5.12 oppure
iOS $\geq$ V1.5.4).

Due precisazioni utili, entrambe verificate:

\begin{itemize}[leftmargin=1.4em,itemsep=2pt]
  \item la funzione ``Detect noise floor'' presente nella schermata dei parametri
        dell'applicazione è \emph{soltanto} una funzione di visualizzazione: mostra i
        valori di rumore misurati ma non li applica come soglie. Non va confusa con
        l'auto-calibrazione;
  \item il documento non descrive alcuna procedura di aggiornamento del firmware, e
        né lo strumento da PC né l'applicazione mobile offrono una funzione di flash.
        La versione del firmware è quindi da considerare un dato dell'esemplare.
\end{itemize}

\chapter{Software realizzato}
\label{app:codice}

Questa appendice documenta l'organizzazione del software prodotto per la tesi e i
comandi con cui è stato usato. I listati completi non sono riportati integralmente:
si indicano struttura, scelte implementative rilevanti e modalità d'uso.

\section{Organizzazione del repository}

\begin{center}\small
\begin{tabularx}{\textwidth}{@{}lX@{}}
\toprule
\textbf{Percorso} & \textbf{Contenuto} \\
\midrule
\texttt{firmware/ld2410b\_logger/} & logger CSV a 5~Hz con engineering mode e
  lettura del PIR --- è il firmware usato per tutta la campagna \\
\texttt{firmware/test00\_ld2410\_info/} & lettura dei parametri del modulo via
  UART: firmware, MAC, soglie, gate, timeout \\
\texttt{firmware/test01\_ld2410\_base/} & sketch diagnostico con auto-scan del baud
  rate e stampa dei byte grezzi \\
\texttt{firmware/test03\_pir\_base/} & verifica del PIR e misura della durata degli
  impulsi \\
\texttt{firmware/test04\_set\_gate/} & configurazione del gate massimo per la prova
  di selettività spaziale \\
\texttt{firmware/ld2420\_monitor/} & lettura del LD2420 in modalità ASCII \\
\texttt{HLK-LD2410x/acquire.py} & acquisizione dalla seriale con metadati
  sperimentali (versione adattata) \\
\texttt{analisi/verifica\_engineering.py} & controllo di qualità dei file CSV \\
\texttt{analisi/analizza\_test.py} & metriche di confronto e aggregazione per
  scenario \\
\texttt{analisi/analizza\_respiro.py} & analisi in frequenza e ricerca del canale
  migliore \\
\bottomrule
\end{tabularx}
\end{center}

\section{Il logger}

\subsection{Scelte implementative rilevanti}

\begin{description}[leftmargin=1.4em,style=nextline,itemsep=3pt]
  \item[Periodo fisso e non ritardo fisso] Il ciclo principale calcola l'istante del
    prossimo campione anziché attendere un intervallo dopo l'elaborazione. È la
    ragione per cui la cadenza misurata è di 200~ms esatti su tutti gli intervalli,
    senza deriva: requisito necessario per l'analisi in frequenza.
  \item[Lettura del PIR nello stesso frame] Il valore del PIR è campionato
    immediatamente prima di comporre la riga CSV, quindi i due sensori condividono la
    base dei tempi per costruzione.
  \item[Engineering mode all'avvio] Il firmware entra in configurazione, attiva
    l'engineering mode, esce dalla configurazione e verifica che i frame ricevuti
    siano del tipo esteso, segnalando l'esito sulla seriale. Un fallimento silenzioso
    produrrebbe file senza colonne per-gate scoperti solo in fase di analisi.
  \item[Intestazione stampata una sola volta] L'intestazione CSV è stampata in
    \texttt{setup()}. Poiché il reset all'apertura della porta non è garantito, lo
    script di acquisizione la ricostruisce dal numero di colonne (si veda oltre).
\end{description}

\subsection{Formato di uscita}

\begin{lstlisting}[caption={Intestazione CSV completa in engineering mode, 29
colonne.},basicstyle=\ttfamily\scriptsize]
timestamp_ms,radar_presence,moving_target,stationary_target,
moving_distance_cm,stationary_distance_cm,moving_energy,stationary_energy,
pir_presence,
menergy_gate0,menergy_gate1,menergy_gate2,menergy_gate3,menergy_gate4,
menergy_gate5,menergy_gate6,menergy_gate7,menergy_gate8,
senergy_gate0,senergy_gate1,senergy_gate2,senergy_gate3,senergy_gate4,
senergy_gate5,senergy_gate6,senergy_gate7,senergy_gate8,
light_level,out_level
\end{lstlisting}

Le prime nove colonne coincidono con quelle del repository di
riferimento~\cite{repoCallisto}, per compatibilità.

\section{Acquisizione}

\subsection{Ambiente}

\begin{lstlisting}[language=bash,caption={Preparazione dell'ambiente Python.},
basicstyle=\ttfamily\scriptsize]
cd HLK-LD2410x
python -m venv .venv
.venv\Scripts\Activate.ps1
pip install pyserial numpy
mkdir data
\end{lstlisting}

\subsection{Uso}

\begin{lstlisting}[language=bash,caption={Esempio di acquisizione di un trial.},
basicstyle=\ttfamily\scriptsize]
python acquire.py --port COM3 --duration 300 ^
  --output data/fermo_seduto_T01.csv ^
  --scenario fermo_seduto --trial T01 ^
  --ground_truth_presence 1 --ground_truth_state still
\end{lstlisting}

Lo script aggiunge a ogni riga le colonne \texttt{pc\_time\_s},
\texttt{group\_id}, \texttt{trial\_id}, \texttt{scenario},
\texttt{ground\_truth\_presence}, \texttt{ground\_truth\_state}.

\subsection{La modifica necessaria}

La versione originale scrive righe solo dopo aver riconosciuto l'intestazione
stampata dall'ESP32 in \texttt{setup()}. Se il reset all'apertura della porta non
scatta, il file risulta \textbf{vuoto senza alcun messaggio}. La versione adattata
ricostruisce l'intestazione dal numero di colonne della prima riga di dati:

\begin{center}\small
\begin{tabular}{@{}rl@{}}
\toprule
\textbf{Colonne} & \textbf{Interpretazione} \\
\midrule
9 & modalità base \\
27 & engineering mode \\
29 & engineering mode con sensore di luce e stato OUT \\
\bottomrule
\end{tabular}
\end{center}

\section{Script di analisi}

\subsection{Controllo di qualità}

\begin{lstlisting}[language=bash,basicstyle=\ttfamily\scriptsize]
python analisi/verifica_engineering.py data/fermo_seduto_T01.csv
\end{lstlisting}

Verifica ed esce con diagnostica su: cadenza effettiva e jitter, presenza delle
colonne per-gate, percentuale di campioni saturi per canale, coerenza fra gate di
picco e distanza riportata, rumore di fondo per-gate nei tratti a stanza vuota,
numero e istante delle transizioni di presenza. Va eseguito su ogni file
\emph{prima} di usarlo per l'analisi.

\subsection{Metriche di confronto}

\begin{lstlisting}[language=bash,basicstyle=\ttfamily\scriptsize]
python analisi/analizza_test.py data/fermo_seduto_T0*.csv --salta-inizio 30
python analisi/analizza_test.py data/ingresso_T0*.csv --event-time 10
python analisi/analizza_test.py data/uscita_T0*.csv --release-time 10
\end{lstlisting}

Le opzioni principali: \texttt{-{}-salta-inizio} scarta i primi secondi (uscita
dell'operatore e stabilizzazione del PIR), \texttt{-{}-event-time} e
\texttt{-{}-release-time} attivano il calcolo delle latenze con la convenzione
dell'evento a istante noto. Con più file dello stesso scenario lo script aggrega e
riporta media e deviazione standard fra trial. Nel calcolo della latenza di rilascio
distingue i casi in cui il rilascio non avviene mai e quelli in cui avviene prima
dell'istante dichiarato, e conta le riaccensioni nella coda.

\subsection{Analisi del respiro}

\begin{lstlisting}[language=bash,basicstyle=\ttfamily\scriptsize]
python analisi/analizza_respiro.py data/fermo_seduto_T01.csv --scan
\end{lstlisting}

La modalità \texttt{-{}-scan} esamina tutti e venti i canali di energia, scarta
quelli saturi e quelli piatti secondo criteri dichiarati, ordina i rimanenti per
rapporto segnale-rumore del picco in banda 0,1--0,5~Hz e riporta la mediana delle
stime concordi con l'elenco dei canali che le sostengono. Esporta lo spettro in CSV
per la produzione dei grafici. Richiede NumPy; gli altri script usano la sola
libreria standard.

\section{Configurazione del gate massimo}

Lo sketch \texttt{firmware/test04\_set\_gate/} imposta il gate massimo via UART per
la prova di selettività spaziale. Tre accorgimenti sono degni di nota, perché
riguardano la scrittura di parametri persistenti su un modulo:

\begin{enumerate}[leftmargin=1.8em,itemsep=2pt]
  \item dopo ogni scrittura \textbf{rilegge sempre} i parametri per conferma, e
        segnala se il valore letto non corrisponde a quello scritto;
  \item \textbf{avvisa} quando la configurazione corrente è diversa da quella di
        fabbrica, per evitare di condurre acquisizioni con parametri modificati
        inconsapevolmente --- il rischio più insidioso in una campagna comparativa;
  \item il comando \texttt{d} riscrive i valori di fabbrica del nostro esemplare
        (le 9+9 soglie di \cref{tab:identita}, gate 8/8 e timeout di 5~s),
        così da poter tornare in ogni momento alla configurazione di riferimento.
\end{enumerate}


\cleardoublepage
\printbibliography[heading=bibintoc,title={Bibliografia e fonti}]

\end{document}
